\chapter{Beams, columns, frames}
\label{vol03:ch08}

\epigraph{It is now apparent in what manner Nature, by means of
greater thickness, can within a moderate length make a strong
support for very great weight; even so, in a giant of unusual size
the bones must either be made of stronger and harder substance than
those of smaller men, or the giant must have stout limbs out of all
proportion to his size.}{Galileo Galilei, \emph{Two New Sciences},
Day Two (1638), in the Drake translation \cite[pp.~127--128]{
hist:galileo1638}.}

\chapmeta{Half-life: durable. Archetypes: stability, scaling. Exercise target: 40. Project track: design plus physical build (hybrid). Hazard class: low (bookshelf-scale loads). Reader path default: standard, with sections explicitly tagged. Named cases: Hyatt Regency walkway collapse (1981, cross-referenced).}

\noindent
A beam carries transverse load by bending; a column carries axial
load that, if compressive and large enough, buckles the column; a
frame combines beams and columns into a structure that carries any
combination of the two. The structural mechanics of these three
elements is the working basis of buildings, bridges, vehicle frames,
furniture, and machine structures. The chapter develops the
theory (Euler-Bernoulli bending, Euler buckling, the slope-deflection
method for indeterminate frames, energy methods), the code-based
discipline that turns the theory into design (load and resistance
factor design, safety factors, deflection limits), and the failure
modes that govern the working envelopes of real structures.

\section{Bending: Euler-Bernoulli theory}\pathtag{core}

A slender beam loaded transversely deforms by bending: the originally
straight neutral axis curves, the fibres above the neutral axis are
compressed, and the fibres below are extended (for a sagging beam).
The \engterm{Euler-Bernoulli beam theory} captures this with three
assumptions:

\begin{itemize}[leftmargin=*]
  \item Plane cross-sections remain plane after deformation.
  \item Cross-sections remain perpendicular to the deformed neutral
    axis.
  \item The material is linear elastic in tension and compression.
\end{itemize}

The first two assumptions, due to Bernoulli, neglect shear
deformation; they are accurate for slender beams (length-to-depth
ratios above about $10$) and become inaccurate for stubby beams,
where the Timoshenko correction (chapter 9 of this volume) is needed.
The third assumption restricts the analysis to stresses below the
material's proportional limit; plastic-bending analysis (chapter 7
of Volume V) handles the post-yield regime.

Under the three assumptions, the strain at distance $y$ from the
neutral axis of a beam with local radius of curvature $\rho$ is
\[
\epsilon = -\frac{y}{\rho},
\]
linear in $y$, with the negative sign reflecting that fibres on the
sagging side (positive $y$ if $y$ points up and the beam sags) are
compressed. Hooke's law gives the stress as $\sigma = E\epsilon =
-Ey/\rho$. The resultant axial force on the cross-section is
$\int \sigma \,dA = -(E/\rho) \int y \,dA = 0$ (by the definition of
the centroid; recall chapter 2 of this volume). The bending moment
on the cross-section is
\[
M = -\int \sigma y \,dA = \frac{E}{\rho}\int y^2 \,dA = \frac{EI}{
\rho},
\]
giving the \engterm{moment-curvature relation}
\[
\frac{1}{\rho} = \kappa = \frac{M}{EI}.
\]
The product $EI$ is the \engterm{flexural rigidity} of the beam, the
product of Young's modulus and the second moment of area about the
centroidal bending axis. The curvature is linearly proportional to
the bending moment; the stiffer the beam (larger $EI$), the smaller
the curvature for a given moment.

Before bending stress can be computed, the bending moment itself must
be known along the span, and that comes from the shear and moment
diagrams of chapter 2 of this volume. Figure
\ref{fig:vol03ch08:shear-moment} shows the canonical case of a central
point load: the shear is piecewise constant and jumps by the load
where it acts, the moment is the running integral of the shear and
peaks where the shear changes sign. Figure
\ref{fig:vol03ch08:udl-shear-moment} shows the uniform-line-load case,
where the shear is linear and the moment parabolic. The script
\path{code/shear_moment.py} integrates $dV/dx=-w$ and $dM/dx=V$ for an
arbitrary mix of point and line loads and writes a sampled profile to
\path{data/shear_moment_profile.csv}; the working engineer reads the
peak moment off such a diagram and carries it into the stress
calculation.

% Figure: shear and bending-moment diagrams for a simply supported beam
% under a central point load P at span L. V is piecewise constant at
% +P/2 and -P/2; M is triangular, peaking at PL/4 at midspan.
\begin{figure}[htbp]
\centering
\begin{tikzpicture}
% loaded beam sketch
\begin{scope}[yshift=4.2cm]
  \draw[engblack, very thick] (0,0) -- (8,0);
  % supports
  \fill[engblue] (0,0) -- (-0.25,-0.45) -- (0.25,-0.45) -- cycle;
  \draw[engblue] (-0.35,-0.6) -- (0.35,-0.6);
  \fill[engblue] (8,0) circle (0.12);
  \draw[engblue] (7.65,-0.45) -- (8.35,-0.45);
  % central load
  \draw[engred, very thick, ->] (4,1.1) -- (4,0.08);
  \node[engred, anchor=south, font=\small] at (4,1.1) {$P$};
  \node[anchor=north, font=\scriptsize] at (0,-0.65) {$R_A=P/2$};
  \node[anchor=north, font=\scriptsize] at (8,-0.6) {$R_B=P/2$};
  \node[anchor=south, font=\scriptsize] at (4,-0.55) {$L/2$};
  \draw[<->, enggray] (0,-0.3) -- (4,-0.3);
\end{scope}
% shear diagram
\begin{scope}[yshift=2.0cm]
  \draw[enggray, ->] (-0.3,0) -- (8.6,0) node[anchor=west, font=\scriptsize]{$x$};
  \node[anchor=east, font=\scriptsize] at (-0.3,0.6) {$V(x)$};
  \draw[engblue, very thick] (0,0.8) -- (4,0.8) -- (4,-0.8) -- (8,-0.8);
  \draw[engblue, thick] (0,0) -- (0,0.8);
  \draw[engblue, thick] (8,-0.8) -- (8,0);
  \node[engblue, anchor=south, font=\scriptsize] at (2,0.8) {$+P/2$};
  \node[engblue, anchor=north, font=\scriptsize] at (6,-0.8) {$-P/2$};
\end{scope}
% moment diagram
\begin{scope}[yshift=-0.2cm]
  \draw[enggray, ->] (-0.3,0) -- (8.6,0) node[anchor=west, font=\scriptsize]{$x$};
  \node[anchor=east, font=\scriptsize] at (-0.3,0.6) {$M(x)$};
  \draw[engred, very thick] (0,0) -- (4,1.2) -- (8,0);
  \fill[engred!12] (0,0) -- (4,1.2) -- (8,0) -- cycle;
  \node[engred, anchor=south, font=\scriptsize] at (4,1.2) {$M_{\max}=PL/4$};
\end{scope}
\end{tikzpicture}
\caption[Shear and moment diagrams, central point load]{Loaded beam,
shear diagram, and bending-moment diagram for a simply supported span
$L$ carrying a central point load $P$. The reactions are each $P/2$.
The shear is constant at $+P/2$ left of the load and $-P/2$ to its
right, jumping by $P$ where the load acts. The moment is the running
integral of the shear, $dM/dx = V$: it rises linearly to a peak of
$PL/4$ at midspan, where the shear changes sign, then falls linearly
to zero at the right support. The peak moment sets the maximum
bending stress $\sigma_{\max}=M_{\max}/S$.}
\label{fig:vol03ch08:shear-moment}
\end{figure}


% Figure: shear and moment diagrams for a simply supported beam under
% a uniform line load w0. Shear is linear from +w0 L/2 to -w0 L/2;
% moment is parabolic, peak w0 L^2/8 at midspan.
\begin{figure}[htbp]
\centering
\begin{tikzpicture}
% loaded beam
\begin{scope}[yshift=4.2cm]
  \draw[engblack, very thick] (0,0) -- (8,0);
  \fill[engblue] (0,0) -- (-0.25,-0.45) -- (0.25,-0.45) -- cycle;
  \fill[engblue] (8,0) circle (0.12);
  \foreach \x in {0.3,1.0,...,7.8}{
    \draw[engred, ->] (\x,0.85) -- (\x,0.08);
  }
  \draw[engred, thick] (0,0.85) -- (8,0.85);
  \node[engred, anchor=south, font=\small] at (4,0.95) {$w_0$ (per unit length)};
\end{scope}
% shear diagram (linear)
\begin{scope}[yshift=2.0cm]
  \draw[enggray, ->] (-0.3,0) -- (8.6,0) node[anchor=west, font=\scriptsize]{$x$};
  \node[anchor=east, font=\scriptsize] at (-0.3,0.6) {$V(x)$};
  \draw[engblue, very thick] (0,0.9) -- (8,-0.9);
  \draw[engblue, thick] (0,0) -- (0,0.9);
  \draw[engblue, thick] (8,-0.9) -- (8,0);
  \node[engblue, anchor=south, font=\scriptsize] at (0.9,0.65) {$+w_0L/2$};
  \node[engblue, anchor=north, font=\scriptsize] at (7.1,-0.65) {$-w_0L/2$};
  \node[enggray, font=\scriptsize, anchor=south] at (4,0.05) {$V=0$ at midspan};
\end{scope}
% moment diagram (parabola)
\begin{scope}[yshift=-0.4cm]
  \draw[enggray, ->] (-0.3,0) -- (8.6,0) node[anchor=west, font=\scriptsize]{$x$};
  \node[anchor=east, font=\scriptsize] at (-0.3,0.7) {$M(x)$};
  \draw[engred, very thick, domain=0:8, samples=60, smooth]
    plot (\x, {1.3*( \x/8 )*(1-\x/8)*4});
  \node[engred, anchor=south, font=\scriptsize] at (4,1.35) {$M_{\max}=w_0L^2/8$};
\end{scope}
\end{tikzpicture}
\caption[Shear and moment, uniform line load]{Simply supported span
$L$ under a uniform line load $w_0$. The shear varies linearly from
$+w_0L/2$ at the left support to $-w_0L/2$ at the right, passing
through zero at midspan. The moment is the integral of the shear and
so is parabolic, peaking where the shear vanishes, with maximum
$w_0L^2/8$ at midspan. The peak moment scales as the square of the
span: doubling $L$ quadruples the design moment for the same line
load.}
\label{fig:vol03ch08:udl-shear-moment}
\end{figure}


The stress at distance $y$ from the neutral axis is then
\[
\sigma = \frac{My}{I},
\]
where the sign convention is now embedded in $M$ (positive for sagging,
negative for hogging) and the formula reads naturally: tension on the
sagging side at positive $y$, compression at negative $y$. The maximum
stress occurs at the extreme fibre, where $y = y_{\max}$, the distance
from the neutral axis to the farthest fibre:
\[
\sigma_{\max} = \frac{M y_{\max}}{I} = \frac{M}{S}, \quad S = \frac{I}{
y_{\max}},
\]
introducing the \engterm{section modulus} $S$, the geometric property
that combines the second moment of area and the extreme-fibre
distance. The section modulus appears throughout structural design
because it is what relates the bending moment to the maximum stress in
one step. Tables of $S$ for standard rolled sections (I-beams, channels,
angles, hollow tubes) appear in engineering handbooks and the AISC
Steel Construction Manual; a short excerpt sufficient for the worked
examples and exercises in this chapter is in
\path{data/standard_sections.csv}, and \path{code/section_properties.py}
computes $A$, $I$, $S$, and the radius of gyration $r$ for rectangular,
circular, hollow, and built-up I sections from first principles.

Figure \ref{fig:vol03ch08:section-properties} makes the geometric point
that governs section choice. The second moment of area rewards both
depth and the placement of material far from the neutral axis. A
rectangle has $I=bh^3/12$; halving the width and doubling the depth
preserves the area yet quadruples $I$, because depth enters cubed. The
rolled I-section concentrates its area in the flanges, where each unit
of material sits at large $y$ and contributes the most to $I$. This is
why beams are deep rather than wide and why efficient sections push
material to the extreme fibres.

% Figure: cross-sections compared by second moment of area. Rectangle,
% same-area I-section, and the depth-cubed scaling. Pure geometry.
\begin{figure}[htbp]
\centering
\begin{tikzpicture}[scale=1.0]
% Rectangle b x h
\begin{scope}[xshift=0cm]
  \draw[engblue, very thick, fill=engblue!10] (0,0) rectangle (1.0,2.0);
  \draw[<->, enggray] (0,-0.3) -- (1.0,-0.3) node[midway, below, font=\scriptsize]{$b$};
  \draw[<->, enggray] (-0.35,0) -- (-0.35,2.0) node[midway, left, font=\scriptsize]{$h$};
  \draw[engred, dashed] (-0.1,1.0) -- (1.1,1.0);
  \node[engblack, anchor=north, font=\scriptsize, align=center] at (0.5,-0.7)
    {rectangle\\ $I=\frac{bh^3}{12}$};
\end{scope}
% Double-depth rectangle
\begin{scope}[xshift=3.0cm]
  \draw[engblue, very thick, fill=engblue!10] (0,-0.5) rectangle (0.5,3.5);
  \draw[<->, enggray] (0,-0.8) -- (0.5,-0.8) node[midway, below, font=\scriptsize]{$b/2$};
  \draw[<->, enggray] (-0.35,-0.5) -- (-0.35,3.5) node[midway, left, font=\scriptsize]{$2h$};
  \node[engblack, anchor=north, font=\scriptsize, align=center] at (0.25,-1.2)
    {same area, deeper\\ $I=\frac{(b/2)(2h)^3}{12}=4\,\frac{bh^3}{12}$};
\end{scope}
% I-section
\begin{scope}[xshift=6.5cm]
  \fill[englime!25] (0,1.7) rectangle (2.4,2.0);   % top flange
  \fill[englime!25] (0,0.0) rectangle (2.4,0.3);   % bottom flange
  \fill[englime!25] (1.0,0.3) rectangle (1.4,1.7); % web
  \draw[engblue, very thick] (0,2.0) -- (2.4,2.0) -- (2.4,1.7) -- (1.4,1.7)
    -- (1.4,0.3) -- (2.4,0.3) -- (2.4,0.0) -- (0,0.0) -- (0,0.3) -- (1.0,0.3)
    -- (1.0,1.7) -- (0,1.7) -- cycle;
  \draw[engred, dashed] (-0.1,1.0) -- (2.5,1.0);
  \node[engblack, anchor=north, font=\scriptsize, align=center] at (1.2,-0.4)
    {I-section\\ material at the\\ extreme fibres};
\end{scope}
\end{tikzpicture}
\caption[Second moment of area and section shape]{The second moment
of area rewards depth and rewards placing material far from the
neutral axis (red dashed). A rectangle has $I=bh^3/12$; halving the
width and doubling the depth keeps the area the same yet quadruples
$I$, because the depth enters cubed. The I-section (right) carries
most of its area in the flanges, far from the neutral axis, where it
contributes most to $I$ per unit of material. This is why structural
members are deep and why rolled sections concentrate steel in the
flanges rather than the web.}
\label{fig:vol03ch08:section-properties}
\end{figure}


The relation between bending moment $M(x)$ and the beam's loading
$w(x)$ was established in chapter 2 of this volume: $dV/dx = -w$ and
$dM/dx = V$, with $w$ the distributed load per unit length, $V$ the
shear, and $M$ the moment. Combined with the moment-curvature relation
and the geometric relation $\kappa \approx d^2 y/dx^2$ for small
deflections (where $y(x)$ is the beam's transverse deflection), the
chain of equations gives
\[
EI \frac{d^4 y}{dx^4} = w(x).
\]
This is the \engterm{Euler-Bernoulli beam equation}: a fourth-order
linear differential equation for the deflection in terms of the load.
The four boundary conditions (two at each end) close the problem.
Common boundary conditions:

\begin{itemize}[leftmargin=*]
  \item Simple support (pin or roller): $y = 0$, $M = EI \,d^2 y/dx^2
    = 0$.
  \item Fixed (clamped, built-in): $y = 0$, $dy/dx = 0$.
  \item Free: $M = EI \,d^2 y/dx^2 = 0$, $V = EI \,d^3 y/dx^3 = 0$.
  \item Guided (roller in a slot): $dy/dx = 0$, $V = 0$.
\end{itemize}

A beam's boundary conditions are part of its mechanical specification;
two beams with the same geometry and load but different end conditions
have different deflections, stresses, and stability.

\begin{principle}[Bending stress is a section-modulus calculation]
For a slender elastic beam under bending moment $M$, the maximum
extreme-fibre stress is $\sigma_{\max} = M/S$, where $S = I/y_{\max}$
is the section modulus. Doubling the depth of a rectangular section
multiplies the section modulus by four; the depth-cubed term in
$I = bh^3/12$ is the structural reason beams are deep rather than wide.
The same calculation, with appropriate sign and orientation, is the
working tool for every bending-dominated structural member in the
book.
\end{principle}

\section{Beam deflection}\pathtag{core}

The Euler-Bernoulli beam equation $EI \,y'''' = w(x)$ can be
integrated four times to obtain $y(x)$ for any loading. The four
constants of integration are set by the four boundary conditions.
For simple loadings and simple boundary conditions, the integral is
straightforward and gives closed-form deflection expressions; for
complex loadings or non-standard boundary conditions, numerical
integration (or the slope-deflection method of section 8.5) is
preferred.

Closed-form deflection results for the standard beam-loading cases,
which the working structural engineer memorises or carries in a
pocket reference:

\begin{itemize}[leftmargin=*]
  \item Simply-supported beam, span $L$, centred point load $P$:
    midspan deflection $\delta = PL^3/(48 EI)$, midspan moment $M =
    PL/4$.
  \item Simply-supported beam, span $L$, uniform line load $w_0$:
    midspan deflection $\delta = 5 w_0 L^4/(384 EI)$, midspan moment
    $M = w_0 L^2/8$.
  \item Cantilever, length $L$, tip point load $P$: tip deflection
    $\delta = PL^3/(3 EI)$, fixed-end moment $M = PL$.
  \item Cantilever, length $L$, uniform line load $w_0$: tip
    deflection $\delta = w_0 L^4/(8 EI)$, fixed-end moment $M = w_0
    L^2/2$.
  \item Fixed-fixed beam, span $L$, centred point load $P$: midspan
    deflection $\delta = PL^3/(192 EI)$, midspan moment $M = PL/8$
    (sagging), end moment $M = PL/8$ (hogging).
  \item Fixed-fixed beam, span $L$, uniform line load $w_0$: midspan
    deflection $\delta = w_0 L^4/(384 EI)$, midspan moment $M =
    w_0 L^2/24$ (sagging), end moment $M = w_0 L^2/12$ (hogging).
\end{itemize}

% Figure: deflection curve y(x) of a simply supported beam under
% uniform load, normalized, against the undeflected axis. Shows the
% 5wL^4/384EI midspan deflection and the symmetric shape.
\begin{figure}[htbp]
\centering
\begin{tikzpicture}
\begin{axis}[
  width=12cm, height=6cm,
  xlabel={position $x/L$}, ylabel={deflection $y/\delta_{\max}$},
  xmin=0, xmax=1, ymin=-1.15, ymax=0.25,
  axis lines=left,
  ytick={0,-0.5,-1}, yticklabels={$0$,$0.5$,$1.0$},
  xtick={0,0.25,0.5,0.75,1},
  tick label style={font=\scriptsize},
  label style={font=\small},
  grid=major, grid style={enggray!18},
]
% undeflected axis
\addplot[enggray, dashed, thick, domain=0:1, samples=2] {0};
% normalized deflection: shape (16/5)(x - 2x^3 + x^4), peak at x=1/2 = 1
\addplot[engblue, very thick, domain=0:1, samples=80, smooth]
  {-( (16/5)*( x - 2*x^3 + x^4) )};
% midspan marker
\addplot[engred, only marks, mark=*, mark options={fill=engred}]
  coordinates {(0.5,-1)};
\node[engred, anchor=west, font=\scriptsize] at (axis cs:0.52,-0.95)
  {$\delta_{\max}=\dfrac{5w_0L^4}{384\,EI}$};
\end{axis}
\end{tikzpicture}
\caption[Deflection curve, uniformly loaded simple beam]{Normalized
deflection shape $y(x)$ of a simply supported beam under uniform line
load, from the closed-form solution
$y(x)=-\dfrac{w_0}{24EI}\,(x^4-2Lx^3+L^3x)$, plotted against the
straight undeflected axis (dashed). The curve is symmetric about
midspan, where the slope is zero and the deflection reaches its
maximum $5w_0L^4/(384EI)$. The fourth-power dependence on span is the
reason long beams sag dramatically: doubling the span multiplies the
midspan deflection by sixteen.}
\label{fig:vol03ch08:deflection-curve}
\end{figure}


Figure \ref{fig:vol03ch08:deflection-curve} plots the deflected shape
for the uniformly loaded simple beam, the workhorse case. The curve is
symmetric, the slope vanishes at midspan, and the maximum deflection is
$5w_0L^4/(384EI)$. The finite-difference solver in
\path{code/beam_deflection.py} reproduces this closed-form midspan
value to better than one part in a thousand on a two-hundred-node grid,
and writes the sampled deflection to \path{data/deflection_profile.csv};
it is the numerical tool to reach for when the loading or the stiffness
varies along the span in ways that defeat a table lookup.

The pattern in the deflection formulas is $\delta \propto L^n / EI$
with $n = 3$ for point loads and $n = 4$ for distributed loads. The
strong dependence on span ($L^3$ or $L^4$) is what makes long beams
deflect dramatically: doubling the span of a uniformly loaded beam
multiplies the deflection by $16$. Increasing the flexural rigidity
$EI$ (by deeper section, better material, or both) is the only
defence; making the load smaller is the only other.

The \engterm{principle of superposition} applies to linear elastic
beams: the deflection under multiple loads is the sum of the
deflections under each load applied separately. The principle is
exact within the linear-elastic, small-deflection assumption set and
is the basis of using tables of single-load deflection results to
analyse complex multi-load problems.

For numerical evaluation of beam deflections without tables, two
standard tools:

\engterm{Moment-area theorems}. Theorem 1: the change in slope
between two points on a beam equals the area of the $M/EI$ diagram
between those two points. Theorem 2: the deflection at one point,
measured from the tangent at another, equals the first moment of the
$M/EI$ diagram between the two points, taken about the point of
measurement. The theorems are useful when only the deflection at
specific points is needed (rather than the full $y(x)$ curve) and
when the loading or stiffness varies along the beam in ways that
make analytical integration tedious.

\engterm{Virtual work for deflections}. The deflection $\delta_P$ at
the point of application of a real load $P$ on a beam equals
$\partial U/\partial P$, where $U$ is the strain energy stored in
the beam (Castigliano's first theorem). For an elastic beam under
bending, $U = \int_0^L M^2/(2EI) \,dx$. The unit-load method
computes the deflection at any point by applying a virtual unit
load at that point and integrating $M(x) m(x)/(EI) \,dx$, where
$M(x)$ is the bending moment from the real load and $m(x)$ from the
virtual unit load. The method extends to combined-loading problems
in section 8.6.

\section{Combined loading}\pathtag{standard}

A real structural member often carries several types of load
simultaneously: axial force, bending in one or two planes, transverse
shear, and torsion. The stress state at any point is the
superposition of the stresses from each load type, by linear
superposition within the elastic range. The superposition produces a
combined-stress state that the engineer must check against the
material's yield criterion (chapter 3 of this volume).

For a beam under axial force $N$, bending moment $M$ about a
centroidal axis, and torsion $T$ about the longitudinal axis (for a
circular shaft), the stresses at a point on the cross-section's
extreme fibre are:

\begin{itemize}[leftmargin=*]
  \item Axial: $\sigma_{\text{axial}} = N/A$, uniform across the
    cross-section.
  \item Bending: $\sigma_{\text{bending}} = M y_{\max}/I = M/S$, at
    the extreme fibre.
  \item Torsion (circular shaft): $\tau_{\text{torsion}} = T r/J$, at
    radius $r$, with $J$ the polar moment of inertia.
  \item Shear (transverse): $\tau_{\text{shear}} = VQ/(Ib)$, with $V$
    the transverse shear force, $Q$ the first moment of area, $I$
    the second moment, and $b$ the width at the location of
    interest.
\end{itemize}

The principal stresses at the extreme fibre then come from the Mohr's
circle of chapter 3 with the local $\sigma_{xx} = \sigma_{\text{axial}}
+ \sigma_{\text{bending}}$, $\sigma_{yy} = 0$, $\sigma_{xy} =
\tau_{\text{torsion}} + \tau_{\text{shear}}$. The von Mises or Tresca
yield criterion (chapter 3) then governs the design.

\engterm{The interaction equation}. Many design codes embed the
combined-loading check in an \engterm{interaction equation} that
relates the loads to a combined utilisation ratio. For axial and
bending,
\[
\frac{N}{N_{\text{allow}}} + \frac{M}{M_{\text{allow}}} \leq 1.0,
\]
with $N_{\text{allow}}$ and $M_{\text{allow}}$ the member's
allowable axial and bending loads from the code. The equation is
linear in the loads; it implicitly assumes elastic behaviour and a
particular ratio of failure modes. More sophisticated codes use
nonlinear interaction equations (squared or cubic terms) to capture
the actual interaction between failure modes for specific member
types. The AISC specification, for instance, uses a bilinear envelope
that steepens when the axial demand exceeds one-fifth of the axial
capacity and flattens below it, the two regimes weighting the two
loads differently according to which one drives the member to its
limit. The envelope and a worked check against it appear in example 8
and figure \ref{fig:vol03ch08:beam-column-interaction}.

\engterm{P-delta effects}. A column in compression that deflects
laterally experiences additional bending moments because the axial
force acts at an eccentricity equal to the lateral deflection. The
additional moment increases the deflection, which increases the
moment, and so on. The amplification factor for the first-order
deflection is approximately $1/(1 - N/N_{\text{cr}})$, where
$N_{\text{cr}}$ is the column's Euler critical load (section 8.4).
The P-delta amplification is small for $N \ll N_{\text{cr}}$ and
catastrophic as $N \to N_{\text{cr}}$. Design codes for slender
columns mandate explicit P-delta analysis when the axial load
exceeds a fraction of the critical load.

\section{Buckling: Euler's formula and beyond}\pathtag{core}

A long slender member in axial compression has a critical load above
which the straight equilibrium is unstable: the member deflects
laterally at constant axial load. The phenomenon is \engterm{buckling},
and its onset is one of the canonical structural-mechanics failure
modes.

For a long elastic column with pinned ends, length $L$, flexural
rigidity $EI$, the \engterm{Euler critical load} is
\[
P_{\text{cr}} = \frac{\pi^2 EI}{L^2}.
\]
The result is derived by writing the moment equilibrium of a slightly
deflected column under axial load $P$, getting the differential
equation $EI \,y'' + P y = 0$, and finding the smallest $P$ for
which a non-zero $y(x)$ satisfies the boundary conditions (Euler,
1744). At $P = P_{\text{cr}}$, the column is in neutral equilibrium:
any small deflection is sustained without growth or decay. At $P >
P_{\text{cr}}$, the deflection grows; the column deflects laterally
until either the deflection produces a bending stress exceeding the
material's yield (yielding-arrested-buckling), or the deflection
amplitude is limited by nonlinear effects (large-deflection theory),
or the column collapses.

For different boundary conditions, the critical load takes the form
\[
P_{\text{cr}} = \frac{\pi^2 EI}{(KL)^2},
\]
with $K$ the \engterm{effective length factor}:

\begin{itemize}[leftmargin=*]
  \item Pinned-pinned: $K = 1.0$.
  \item Fixed-fixed: $K = 0.5$ (twice the elastic stiffness, four
    times the critical load).
  \item Fixed-free (cantilever): $K = 2.0$ (quarter the critical
    load).
  \item Fixed-pinned: $K = 0.7$.
\end{itemize}

% Figure: four classic column end conditions with buckled shapes and
% effective-length factors K. Pinned-pinned, fixed-fixed, fixed-free,
% fixed-pinned. Buckled shapes drawn as sine-like curves.
\begin{figure}[htbp]
\centering
\begin{tikzpicture}[scale=1.0]
% Pinned-pinned, K=1.0 : half sine over full length
\begin{scope}[xshift=0cm]
  \draw[enggray, dashed] (0,0) -- (0,4);
  \draw[engblue, very thick, domain=0:4, samples=40, smooth]
    plot ({0.6*sin(deg(pi*\x/4))}, \x);
  \fill[engblack] (0,0) circle (0.08);
  \fill[engblack] (0,4) circle (0.08);
  \draw[engblack] (-0.3,-0.15) -- (0.3,-0.15);
  \node[anchor=north, font=\scriptsize, align=center] at (0,-0.4)
    {pinned-pinned\\ $K=1.0$};
\end{scope}
% Fixed-fixed, K=0.5 : full sine, inflection points at L/4 and 3L/4
\begin{scope}[xshift=3.2cm]
  \draw[enggray, dashed] (0,0) -- (0,4);
  \draw[engred, very thick, domain=0:4, samples=40, smooth]
    plot ({0.5*(1-cos(deg(2*pi*\x/4)))*0.5*( \x<2 ? 1 : -1)}, \x);
  \draw[engblack, line width=2pt] (-0.35,0) -- (0.35,0);
  \draw[engblack, line width=2pt] (-0.35,4) -- (0.35,4);
  \node[anchor=north, font=\scriptsize, align=center] at (0,-0.4)
    {fixed-fixed\\ $K=0.5$};
\end{scope}
% Fixed-free (cantilever), K=2.0 : quarter sine
\begin{scope}[xshift=6.4cm]
  \draw[enggray, dashed] (0,0) -- (0,4);
  \draw[engamber, very thick, domain=0:4, samples=40, smooth]
    plot ({0.7*(1-cos(deg(pi*\x/8)))}, \x);
  \draw[engblack, line width=2pt] (-0.35,0) -- (0.35,0);
  \node[anchor=north, font=\scriptsize, align=center] at (0,-0.4)
    {fixed-free\\ $K=2.0$};
\end{scope}
% Fixed-pinned, K=0.7
\begin{scope}[xshift=9.6cm]
  \draw[enggray, dashed] (0,0) -- (0,4);
  \draw[englime!70!engblack, very thick, domain=0:4, samples=40, smooth]
    plot ({0.55*sin(deg(1.43*pi*\x/4))*(1-\x/5)}, \x);
  \draw[engblack, line width=2pt] (-0.35,0) -- (0.35,0);
  \fill[engblack] (0,4) circle (0.08);
  \draw[engblack] (-0.3,4.18) -- (0.3,4.18);
  \node[anchor=north, font=\scriptsize, align=center] at (0,-0.4)
    {fixed-pinned\\ $K=0.7$};
\end{scope}
\end{tikzpicture}
\caption[Column end conditions and effective length]{The four
classical column end conditions with their buckled mode shapes
(exaggerated) and effective-length factors $K$. The critical load is
$P_{\mathrm{cr}}=\pi^2EI/(KL)^2$, so the capacity scales as $1/K^2$.
The pinned-pinned column buckles in a half sine with $K=1.0$; the
fixed-fixed column has inflection points at the quarter and
three-quarter points, halving the effective length to $K=0.5$ and
quadrupling the critical load; the cantilever (fixed-free) has
$K=2.0$ and only a quarter of the pinned-pinned capacity; the
fixed-pinned column sits between, at $K=0.7$. End restraint is the
single most powerful design parameter for a slender column.}
\label{fig:vol03ch08:column-end-conditions}
\end{figure}


Figure \ref{fig:vol03ch08:column-end-conditions} draws the four buckled
mode shapes with their effective-length factors. Because the critical
load scales as $1/K^2$, end restraint is the most powerful design
parameter for a slender column: fixing both ends quadruples the
capacity of an otherwise identical pinned-pinned column, while leaving
one end free as a flagpole cantilever cuts it to a quarter. The script
\path{code/euler_buckling.py} computes $P_{\text{cr}}$ for all four
cases and the transition slenderness, and tabulates the strength curve.

The effective length $KL$ is the distance between adjacent inflection
points of the buckled shape. A doubly-pinned column buckles in a
half-sine, $KL = L$; a fixed-fixed column has inflection points at
quarter and three-quarter span, $KL = L/2$. The effective-length
factor is the central quantity in column-design tables and is
extended in design codes to account for end-restraint conditions
intermediate between the idealised cases.

\engterm{Slenderness ratio}. The dimensionless
\[
\lambda = \frac{KL}{r}, \quad r = \sqrt{I/A},
\]
with $r$ the section's radius of gyration, governs which failure
mode (elastic buckling or material yielding) governs. The transition
slenderness $\lambda_c$ at which the critical buckling stress equals
the yield stress is
\[
\lambda_c = \pi\sqrt{\frac{E}{\sigma_y}}.
\]
For mild structural steel ($E = 200\,\text{GPa}$, $\sigma_y =
250\,\text{MPa}$), $\lambda_c \approx 89$.

% Figure: column strength curve. Critical stress vs slenderness ratio.
% Euler hyperbola sigma_cr = pi^2 E / lambda^2, capped by yield, with
% the imperfection-reduced design curve below both.
\begin{figure}[htbp]
\centering
\begin{tikzpicture}
\begin{axis}[
  width=12cm, height=7.5cm,
  xlabel={slenderness ratio $\lambda=KL/r$},
  ylabel={critical stress $\sigma_{\mathrm{cr}}$ (MPa)},
  xmin=0, xmax=200, ymin=0, ymax=300,
  axis lines=left,
  legend pos=north east, legend cell align=left,
  legend style={font=\scriptsize},
  tick label style={font=\scriptsize}, label style={font=\small},
  grid=major, grid style={enggray!18},
]
% yield cap at 250 MPa for short columns
\addplot[engred, very thick, domain=0:89, samples=2] {250};
\addlegendentry{yield cap $\sigma_y=250$}
% Euler hyperbola: pi^2 * 200000 / lambda^2  (E in MPa)
\addplot[engblue, very thick, domain=89:200, samples=80]
  {pi^2*200000/x^2};
\addlegendentry{Euler $\pi^2E/\lambda^2$}
% transition slenderness marker lambda_c ~ 89
\addplot[enggray, dashed, thick] coordinates {(89,0) (89,250)};
\node[enggray, font=\scriptsize, anchor=west] at (axis cs:91,40)
  {$\lambda_c=\pi\sqrt{E/\sigma_y}\approx 89$};
% imperfection-reduced design curve (illustrative, below both)
\addplot[engamber, very thick, dashed, domain=0:200, samples=80]
  {min(250, pi^2*200000/(x^2))*0.78};
\addlegendentry{design curve (imperfections)}
\node[engblack, font=\scriptsize, anchor=south] at (axis cs:45,255) {short: yielding};
\node[engblack, font=\scriptsize, anchor=west] at (axis cs:130,90) {long: Euler buckling};
\end{axis}
\end{tikzpicture}
\caption[Column strength versus slenderness]{Critical stress as a
function of slenderness ratio. Short columns ($\lambda<\lambda_c$)
are capped by material yielding at $\sigma_y$; long columns
($\lambda>\lambda_c$) follow the Euler hyperbola
$\sigma_{\mathrm{cr}}=\pi^2E/\lambda^2$. The transition slenderness
$\lambda_c=\pi\sqrt{E/\sigma_y}\approx 89$ for mild steel marks where
the two curves cross. Real columns (amber dashed) carry less than
either ideal because of initial out-of-straightness, load
eccentricity, and residual stresses; the reduction is largest in the
intermediate range, near $\lambda_c$, where neither pure yielding nor
pure elastic buckling dominates. Design codes calibrate the lower
curve to test data.}
\label{fig:vol03ch08:column-curve}
\end{figure}


Figure \ref{fig:vol03ch08:column-curve} plots the column strength
curve: the Euler hyperbola $\sigma_{\text{cr}}=\pi^2E/\lambda^2$ for
long columns, the yield cap for short ones, and the imperfection-reduced
design curve that real columns follow, written to
\path{data/column_strength_curve.csv} by the buckling script. Columns
with $\lambda >
\lambda_c$ are \engterm{long columns}, governed by elastic Euler
buckling. Columns with $\lambda < \lambda_c$ are \engterm{short
columns}, governed by yielding (sometimes with imperfection-
sensitive failure modes that produce lower-than-yield strengths).
The transition region is the \engterm{intermediate-column} regime,
empirically fitted with the Johnson formula or the Perry-Robertson
formula in design codes.

\engterm{Imperfections and the real column}. The Euler analysis
assumes a perfectly straight initial geometry, a perfectly axial
load, and a perfectly linear-elastic material. Real columns have
initial out-of-straightness (typically $L/1000$ in well-rolled steel
sections), load eccentricities, residual stresses from rolling, and
nonlinear material behaviour near yield. Each of these reduces the
column's actual buckling load below the Euler prediction; the
reduction is approximately $30\,\%$ for typical structural steel
columns in the intermediate slenderness range. Design codes account
for the imperfection-driven reduction through column curves (AISC,
Eurocode 3) calibrated to extensive test data.

\engterm{Lateral-torsional buckling}. A deep narrow beam under
bending can buckle out of its plane of loading by combined lateral
displacement and twist. The mechanism is analogous to column
buckling but in a more complex geometry. Lateral-torsional buckling
governs the design of long unbraced beams; the design check is
either a closed-form critical-moment formula or an interaction with
material yielding through code curves. The deeper the beam, the more
prone it is to lateral-torsional buckling and the more important
lateral bracing becomes.

% Figure: lateral-torsional buckling of a deep narrow beam. In-plane
% bending vs out-of-plane lateral displacement plus twist. Two small
% 3D-ish sketches: stable (loaded in plane) and buckled (lateral+twist).
\begin{figure}[htbp]
\centering
\begin{tikzpicture}[scale=1.0]
% Stable: deep narrow beam loaded in plane
\begin{scope}[xshift=0cm]
  \draw[engblue, very thick] (0,0) -- (4,0);          % span
  \draw[engblue, thick] (0,-0.05) -- (0,0.9);          % left depth
  \draw[engblue, thick] (4,-0.05) -- (4,0.9);          % right depth
  \draw[engblue, very thick] (0,0.9) -- (4,0.9);       % top fibre
  \draw[engred, very thick, ->] (2,1.7) -- (2,0.95);
  \node[engred, anchor=south, font=\small] at (2,1.7) {$M$};
  \fill[engblack] (0,0) circle (0.07);
  \fill[engblack] (4,0) circle (0.07);
  \node[anchor=north, font=\scriptsize, align=center] at (2,-0.4)
    {below $M_{\mathrm{cr}}$:\\ bends in plane};
\end{scope}
% Buckled: lateral displacement + twist
\begin{scope}[xshift=6cm]
  % perspective span sweeping sideways at midspan
  \draw[engamber, very thick] (0,0)
    .. controls (1.6,0.0) and (2.4,0.55) .. (4,0);
  \draw[engamber, thick] (0,0.9)
    .. controls (1.6,0.9) and (2.4,1.0) .. (4,0.9);
  % depth lines, twisted at midspan
  \draw[engamber, thick] (0,0) -- (0,0.9);
  \draw[engamber, thick] (4,0) -- (4,0.9);
  \draw[engamber, thick] (2.0,0.27) -- (2.4,1.0); % twisted mid section
  \draw[engred, very thick, ->] (2,1.8) -- (2,1.05);
  \node[engred, anchor=south, font=\small] at (2,1.8) {$M_{\mathrm{cr}}$};
  % lateral displacement arrow
  \draw[engblue, ->] (2.0,-0.55) -- (2.4,-0.05);
  \node[engblue, font=\scriptsize, anchor=north] at (2.0,-0.6) {lateral + twist};
  \fill[engblack] (0,0) circle (0.07);
  \fill[engblack] (4,0) circle (0.07);
  \node[anchor=north, font=\scriptsize, align=center] at (2,-1.0)
    {at $M_{\mathrm{cr}}$:\\ buckles sideways};
\end{scope}
\end{tikzpicture}
\caption[Lateral-torsional buckling of a deep beam]{A deep narrow
beam bent about its strong axis is stiff in plane but soft against
sideways movement. Below the critical moment $M_{\mathrm{cr}}$ it
bends in its plane of loading (left). At $M_{\mathrm{cr}}$ the
compression flange, acting like a column with the tension flange and
web restraining it, buckles out of plane: the beam displaces
laterally and twists at once (right). The deeper and narrower the
section, and the longer the unbraced length, the lower
$M_{\mathrm{cr}}$. Lateral bracing of the compression flange shortens
the effective unbraced length and is the standard defence; this is
why floor beams are braced by the deck they carry.}
\label{fig:vol03ch08:lateral-torsional}
\end{figure}


Figure \ref{fig:vol03ch08:lateral-torsional} shows the mechanism. The
compression flange of a beam bent about its strong axis behaves like a
column held by the web and tension flange; when the moment reaches the
critical value the flange buckles sideways, dragging the section into a
combined lateral displacement and twist. The defence is to brace the
compression flange at intervals short enough to keep the unbraced
length below the limit at which lateral-torsional buckling governs. A
concrete floor slab cast onto the top flange of a simply supported beam
braces it continuously, which is why composite floor beams rarely fail
this way in service yet can buckle during construction before the slab
is poured.

\begin{principle}[Slenderness limits compression capacity]
A compression member's load capacity depends on its slenderness
ratio $\lambda = KL/r$. Short members fail by material yielding;
long members fail by elastic buckling at the Euler load $P_{cr}
= \pi^2 EI/(KL)^2$; intermediate members are governed by an
interaction of yielding and buckling. Designing a compression
member is not stress-design; it is stability-design plus
imperfection allowance plus material yielding, with the dominant
mode set by the slenderness ratio.
\end{principle}

\section{Frames: indeterminate structures}\pathtag{standard}

A frame is an assembly of beams and columns connected at rigid or
semi-rigid joints. Frames in general are statically indeterminate:
the global equilibrium equations are fewer than the unknown
reactions and internal forces; the additional equations come from
compatibility of deformations.

The standard analysis methods:

\engterm{Slope-deflection method}. For a frame with rigid joints,
each member's end moments are expressed in terms of the rotations
and translations of the joints at the member's ends. The
slope-deflection equation for a prismatic member $AB$ writes the moment
at end $A$ as
\[
M_{AB} = \frac{2EI}{L}\left(2\theta_A + \theta_B - 3\psi\right)
+ M^{F}_{AB},
\]
where $\theta_A$ and $\theta_B$ are the rotations of the two joints,
$\psi = \Delta/L$ is the chord rotation produced by a relative
transverse displacement $\Delta$ of the ends (the sidesway term), and
$M^{F}_{AB}$ is the fixed-end moment from the loads applied along the
member. The companion equation for end $B$ swaps the roles of $\theta_A$
and $\theta_B$. Compatibility at
each joint (the rotations of all members meeting at a joint are the
same as the joint's rotation) and equilibrium at each joint (the
sum of member end-moments equals the applied joint moment) produce
enough equations to solve. For a frame free to sway, one extra
equilibrium equation per independent sway degree of freedom, the
horizontal shear balance of a storey, closes the system. The unknowns
are the joint rotations and the sway displacements; once they are
found, back-substitution into the slope-deflection equations gives every
member end moment. The method works analytically for small
frames and is the algebraic ancestor of the stiffness method (below)
for larger frames, which packages the same end-moment relations into
matrix form.

\engterm{Stiffness method}. The frame is discretised into members
connected at joints; each member has a $4 \times 4$ or $6 \times 6$
stiffness matrix relating its end forces and moments to its end
displacements and rotations. The global stiffness matrix assembles
from the member matrices by accounting for the geometric
relationship between member and global coordinate systems; the
global equilibrium equations $K \vec{u} = \vec{F}$ relate the
unknown joint displacements to the applied joint forces. The method
is the foundation of finite-element structural analysis and is
implemented in every structural-analysis software package.

\engterm{Force method (flexibility method)}. The dual of the
stiffness method: a statically determinate primary structure is
chosen by removing enough redundant reactions, the deflections of the
primary structure under the applied loads are computed, the
redundant reactions are reapplied as unknowns, and the compatibility
conditions at the released positions give equations for the unknowns.
The force method is illuminating for hand analysis of indeterminate
structures with few redundancies; it is rarely the choice for
computer analysis.

For a frame of $n$ joints in two dimensions, the global stiffness
matrix is $3n \times 3n$ before boundary conditions are imposed (two
translations and one rotation per joint); after imposing supports
the system is solvable for the remaining unknowns. Frames with
hundreds of joints are routine in commercial structural-analysis
software; frames with millions of joints are routine in finite-
element analysis of complex structures.

% Figure: portal frame under lateral load, undeformed and sway-deformed
% shapes, with the portal-method assumption of inflection points at
% mid-height of columns and mid-span of the beam.
\begin{figure}[htbp]
\centering
\begin{tikzpicture}[scale=1.0]
% undeformed frame
\draw[engblack, very thick] (0,0) -- (0,3.4);   % left column
\draw[engblack, very thick] (5,0) -- (5,3.4);   % right column
\draw[engblack, very thick] (0,3.4) -- (5,3.4); % beam
% fixed bases
\draw[engblack, line width=2pt] (-0.4,0) -- (0.4,0);
\draw[engblack, line width=2pt] (4.6,0) -- (5.4,0);
\foreach \x in {-0.4,-0.15,0.1,0.35}{ \draw[enggray] (\x,0) -- (\x-0.2,-0.25); }
\foreach \x in {4.6,4.85,5.1,5.35}{ \draw[enggray] (\x,0) -- (\x-0.2,-0.25); }
% lateral load
\draw[engred, very thick, ->] (-1.2,3.4) -- (-0.08,3.4);
\node[engred, anchor=east, font=\small] at (-1.2,3.4) {$H$};
% swayed (deformed) shape, dashed, exaggerated
\draw[engblue, thick, dashed] (0,0) .. controls (0.35,1.7) .. (0.55,3.4);
\draw[engblue, thick, dashed] (5,0) .. controls (5.35,1.7) .. (5.55,3.4);
\draw[engblue, thick, dashed] (0.55,3.4) -- (5.55,3.4);
% inflection-point markers (portal method)
\fill[engamber] (0,1.7) circle (0.09);
\fill[engamber] (5,1.7) circle (0.09);
\fill[engamber] (2.5,3.4) circle (0.09);
\node[engamber, font=\scriptsize, anchor=west] at (0.12,1.7) {inflection};
\node[engamber, font=\scriptsize, anchor=south] at (2.5,3.5) {inflection};
% dimensions
\draw[<->, enggray] (0,-0.6) -- (5,-0.6) node[midway, below, font=\scriptsize]{span $L$};
\draw[<->, enggray] (6.0,0) -- (6.0,3.4) node[midway, right, font=\scriptsize]{height $h$};
\end{tikzpicture}
\caption[Portal frame sway under lateral load]{A fixed-base portal
frame under a lateral load $H$. The undeformed frame (solid) sways
sideways into the deflected shape (blue dashed). The portal method of
approximate analysis places inflection points (amber) at the
mid-height of each column and the mid-span of the beam, where the
bending moment passes through zero, and assumes the horizontal shear
divides between the columns in proportion to their stiffness. These
assumptions turn an indeterminate frame into a determinate one that
can be solved by statics alone, accurately enough for preliminary
low-rise design.}
\label{fig:vol03ch08:portal-frame}
\end{figure}


\engterm{Approximate analysis methods}. Before computer analysis,
several approximate hand-analysis methods were standard: the
portal frame method (for low-rise frames under lateral load), the
cantilever method (for tall frames under lateral load), and the
moment-distribution method (Hardy Cross's iterative method, 1932,
which gives the moments to any desired accuracy by successive
approximation). The moment-distribution method is no longer the
working tool but remains a check method and a useful conceptual
framework for understanding load paths in indeterminate structures.
Figure \ref{fig:vol03ch08:portal-frame} shows the portal method's
central assumption: inflection points (zero moment) at the mid-height
of the columns and the mid-span of the beam, which reduce an
indeterminate frame to a determinate one solvable by statics alone.

\begin{mastery}
\textbf{The direct stiffness method, member by member.} The
computer-implemented stiffness method is worth seeing once in full,
because every structural-analysis package is this calculation at scale.
Each 2D frame member has three degrees of freedom at each end: two
translations and one rotation, six in all. In the member's own local
axes, the $6\times 6$ stiffness matrix couples axial action (the
$EA/L$ terms) to flexural action (the $12EI/L^3$, $6EI/L^2$, $4EI/L$,
$2EI/L$ terms familiar from slope-deflection):
\[
k_{\text{local}} =
\begin{pmatrix}
\frac{EA}{L} & 0 & 0 & -\frac{EA}{L} & 0 & 0 \\[2pt]
0 & \frac{12EI}{L^3} & \frac{6EI}{L^2} & 0 & -\frac{12EI}{L^3} & \frac{6EI}{L^2} \\[2pt]
0 & \frac{6EI}{L^2} & \frac{4EI}{L} & 0 & -\frac{6EI}{L^2} & \frac{2EI}{L} \\[2pt]
-\frac{EA}{L} & 0 & 0 & \frac{EA}{L} & 0 & 0 \\[2pt]
0 & -\frac{12EI}{L^3} & -\frac{6EI}{L^2} & 0 & \frac{12EI}{L^3} & -\frac{6EI}{L^2} \\[2pt]
0 & \frac{6EI}{L^2} & \frac{2EI}{L} & 0 & -\frac{6EI}{L^2} & \frac{4EI}{L}
\end{pmatrix}.
\]
A rotation matrix $T$ built from the member's direction cosines
$c=\cos\theta$, $s=\sin\theta$ transforms this into global axes,
$k_{\text{global}} = T^{\!\top} k_{\text{local}} T$. The member
matrices are then added into a global $3n\times 3n$ matrix $K$ by
placing each member's entries at the rows and columns of the two
joints it connects (assembly by the connectivity). Supports are
imposed by deleting the rows and columns of the restrained degrees of
freedom; what remains is solved as $K\vec{u}=\vec{F}$ for the unknown
joint displacements, and the member end forces follow from
$k_{\text{global}}\vec{u}_{\text{member}}$. The implementation in
\path{code/frame_stiffness.py} carries out exactly this sequence for a
single member and checks the joint rotation against the closed-form
value for a propped cantilever; extending it to many members is
bookkeeping, not new physics.
\end{mastery}

\section{Energy methods for structures}\pathtag{standard}

Chapter 5 of this volume introduced energy methods in the context of
analytical mechanics. The same methods apply to structural mechanics
with specific structural form.

\engterm{Strain energy in bending}. A beam of length $L$ under
bending moment $M(x)$ stores strain energy
\[
U = \int_0^L \frac{M^2(x)}{2 EI} \,dx.
\]
For other load types the corresponding strain energy is $\int N^2/(2
EA) \,dx$ for axial, $\int T^2/(2GJ) \,dx$ for torsion, and $\int
fV^2/(2 GA) \,dx$ for shear (with $f$ a shape factor for the
cross-section, typically $\sim 6/5$ for a rectangular section). For
combined loading the strain energies sum.

\engterm{Castigliano's first theorem}. For an elastic structure with
applied loads $P_1, P_2, \ldots, P_n$ producing displacements $\delta_1,
\delta_2, \ldots, \delta_n$ at the load points, the displacement
$\delta_i$ at the $i$th load equals the partial derivative of the
total strain energy with respect to $P_i$:
\[
\delta_i = \frac{\partial U}{\partial P_i}.
\]
The theorem is a direct application of the principle of virtual work
for structures.

\engterm{Castigliano's second theorem (least work)}. For a statically
indeterminate elastic structure, the redundant reactions are those
values that minimise the strain energy with respect to the
redundants. The theorem provides a variational route to
indeterminate analysis: choose the redundants as unknowns,
differentiate the strain energy with respect to each redundant, set
the derivative to zero, and solve. The method gives the same
answer as the force method but with the discipline of variational
calculation.

\engterm{The unit-load method}. To find the deflection at any point
on an elastic structure under a real loading, apply a virtual unit
load at that point in the direction of the desired deflection, compute
the resulting internal-force distribution (from statics, no
deformation analysis required), and integrate
\[
\delta = \int_0^L \frac{M_{\text{real}}(x) \cdot m_{\text{virtual}}(x)
}{EI} \,dx + \text{terms for axial, torsion, shear},
\]
with $M_{\text{real}}$ from the real load and $m_{\text{virtual}}$
from the unit load. The unit-load method is the standard hand
analysis tool for beam and frame deflections; it generalises to
truss displacements (with the axial-energy term only) and to
combined-loading structures.

The energy methods are mathematically equivalent to the
moment-curvature integration of section 8.2, but they handle
combined loading, indeterminate structures, and complex geometry
more cleanly. They are the standard tool for hand analysis when a
structure has more than two or three members; for larger structures,
the stiffness method (computer-implemented) is the working tool.

\begin{mastery}
\textbf{Plastic hinges and limit analysis.} The elastic analysis of
the preceding sections stops at first yield, the moment $M_y=\sigma_y
S$ at which the extreme fibre reaches the yield stress. A ductile steel
beam carries more load past that point: yielding spreads inward from
the extreme fibres until the whole section is at the yield stress, at
which the section forms a \engterm{plastic hinge} carrying the plastic
moment
\[
M_p = \sigma_y Z, \qquad Z = \int_A |y|\,dA,
\]
where $Z$ is the plastic section modulus, the first moment of the area
about the equal-area axis. For a rectangle $Z=bh^2/4$ against the
elastic $S=bh^2/6$, so $M_p/M_y = Z/S = 1.5$: a rectangular beam
carries fifty per cent more moment after first yield before the hinge
forms. This ratio is the \engterm{shape factor}; it is about $1.1$ to
$1.2$ for rolled I-sections, whose material already sits near the
extreme fibres.

A statically determinate beam collapses when one plastic hinge forms.
An indeterminate beam survives the first hinge, which simply removes
one redundancy and redistributes moment to the still-elastic regions;
collapse comes only when enough hinges form to turn the structure into
a mechanism. A fixed-fixed beam under a central point load, for
instance, forms hinges first at the two fixed ends (the high-moment
locations) and finally at midspan, three hinges making a mechanism. The
\engterm{upper-bound (kinematic) theorem} of plastic limit analysis
equates the external work done by the collapse load moving through a
virtual mechanism to the internal work dissipated at the hinges,
$\sum M_p\,\theta_i$, and the lowest collapse load over all admissible
mechanisms is the true one. Limit-state design codes use the plastic
moment, not the first-yield moment, for the strength of ductile
members, which is the deeper reason that ductility is prized in
structural materials: a ductile structure warns and redistributes
before it collapses, while a brittle one fails at first crack.
\end{mastery}

\section{Code-based design preview}\pathtag{mastery}

The previous sections give the structural mechanics for analysing
loads, stresses, deflections, and stability. The conversion of that
analysis into a design that can be built and approved follows the
discipline of \engterm{building codes}: national or jurisdictional
standards that specify how loads are determined, what safety factors
apply, and what allowable stresses, deflections, and slenderness
ratios are. The code's discipline is regulatory, not physical: it
encodes the engineering profession's accumulated experience of what
designs are safe enough to approve, with the experience including
many failures whose lessons the code now embodies.

\engterm{Load and resistance factor design (LRFD)}. The dominant
modern format. Design loads are factored upward (e.g., dead load
multiplied by $1.2$, live load by $1.6$) to account for load
uncertainty; member capacities are factored downward (multiplied by
a resistance factor $\phi < 1$, typically $0.9$ for tension, $0.75$
to $0.9$ for compression depending on slenderness) to account for
material and geometric variability. The design check is
\[
\sum (\gamma_i Q_i) \leq \phi R_n,
\]
factored demand on the left, factored capacity on the right. The
inequality must be satisfied at every load combination the code
requires.

\engterm{Allowable stress design (ASD)}. The older format, still
used in some jurisdictions and for some materials, wood and
masonry among them. Material strength is divided by a single safety factor
(typically $1.67$ to $2.0$ for steel) to give the allowable stress;
the design check is that the working stress under unfactored loads
not exceed the allowable. ASD is conceptually simpler but
statistically less sound than LRFD; it cannot easily distinguish
between high-variability and low-variability loads.

\engterm{Deflection limits}. Beyond strength, codes limit the
deflection of structural members to limit damage to non-structural
elements (interior partitions, glazing, ceilings), occupant
perception, and dynamic effects. Typical limits: $L/360$ for
residential floor live load (a $4\,\si{\meter}$ span beam may
deflect no more than $11\,\si{\milli\meter}$ under live load alone);
$L/240$ for total dead-plus-live load; $L/500$ or $L/600$ for
glazing-supporting beams. The limits are working numbers, not
fundamental, and the appropriate limit depends on the supported
elements.

\engterm{Code references}. Standard codes the reader will encounter:

\begin{itemize}[leftmargin=*]
  \item AISC 360 (Specification for Structural Steel Buildings),
    United States.
  \item Eurocode 3 (EN 1993, Design of Steel Structures), European
    Union.
  \item ACI 318 (Building Code Requirements for Structural Concrete),
    United States.
  \item Eurocode 2 (EN 1992, Design of Concrete Structures),
    European Union.
  \item NDS (National Design Specification for Wood Construction),
    United States.
  \item Eurocode 5 (EN 1995, Design of Timber Structures), European
    Union.
\end{itemize}

The code's discipline is regulatory, not aspirational: a design that
satisfies the code is approved; a design that does not is not. The
working engineer carries the code and reads it in the way the
working chemist reads the periodic table.

\begin{estimation}
\textbf{Estimate the maximum span of a $50\,\si{\milli\meter} \times
200\,\si{\milli\meter}$ pine floor joist carrying a residential
live load of $2\,\text{kN}/\text{m}$.}

\smallskip
\textit{Estimate, before reading on.}

\smallskip
For a simply supported beam with uniform load $w_0$, the midspan
bending moment is $w_0 L^2 / 8$ and the midspan deflection is
$5 w_0 L^4 / (384 EI)$. The two design criteria are typically
bending stress and deflection.

\textbf{Bending criterion.} For pine, the allowable bending stress
is approximately $\sigma_{\text{allow}} \approx 10\,\text{MPa}$
(varies by species and grade). The maximum bending stress is
$\sigma = M/S$, where for a rectangular section $S = bh^2/6$. With
$b = 50\,\si{\milli\meter}$ and $h = 200\,\si{\milli\meter}$,
$S = 50 \times 200^2/6 \approx 3.33 \times 10^5\,\si{\milli\meter
\cubed} = 3.33 \times 10^{-4}\,\si{\meter\cubed}$. Setting
$\sigma_{\text{allow}} = w_0 L^2 / (8 S)$:
\[
L^2 = \frac{8 S \sigma_{\text{allow}}}{w_0} = \frac{8 \times 3.33
\times 10^{-4} \times 10^7}{2000} = 13.3\,\si{\meter\squared},
\]
giving $L \approx 3.65\,\si{\meter}$.

\textbf{Deflection criterion.} A typical residential deflection limit
is $L/360$ under live load. With $E = 10\,\text{GPa}$ for pine and
$I = bh^3/12 = 50 \times 200^3/12 \approx 3.33 \times 10^7\,\si{
\milli\meter\tothe{4}} = 3.33 \times 10^{-5}\,\si{\meter\tothe{4}}$:
\[
\frac{L}{360} = \frac{5 w_0 L^4}{384 EI}, \quad L^3 = \frac{384 EI}{
1800 w_0} = \frac{384 \times 10^{10} \times 3.33 \times 10^{-5}}{
1800 \times 2000} \approx 35.5\,\si{\meter\cubed},
\]
giving $L \approx 3.29\,\si{\meter}$.

The deflection criterion governs at $L \approx 3.3\,\si{\meter}$
rather than the bending criterion at $L \approx 3.65\,\si{\meter}$.
The working design span is the smaller of the two, with a margin
for the load uncertainty.

The postmortem: the deflection criterion frequently governs
residential floor joist design at modest spans, because human
perception of floor "bounce" is sensitive to deflection at frequencies
in the walking band. Stiffer beams (deeper section, higher $E$
material) defer the deflection limit; the cube on depth ($I \propto
h^3$) makes depth the most powerful parameter. Doubling the depth
from $200$ to $400\,\si{\milli\meter}$ multiplies $I$ (and the
deflection-limited span cubed) by eight, giving a deflection-limited
span of $\sim 6.6\,\si{\meter}$.
\end{estimation}

\section{Worked examples}\pathtag{standard}

The methods of the chapter come together in full when the bending,
deflection, buckling, and combined-load checks are carried out
end to end on concrete numbers. The five examples below are worked at
the level of detail a design check requires, with every substitution
shown.

\paragraph{Example 1: a simply supported floor beam, two criteria.}
A simply supported steel beam spans $L=6\,\si{\meter}$ and carries a
uniform service load $w_0 = 12\,\text{kN}/\text{m}$ (dead plus live).
Take $E = 200\,\text{GPa}$ and a deflection limit of $L/360$. Find the
required section modulus and the required second moment of area, then
pick a section from \path{data/standard_sections.csv}.

The maximum moment is
\[
M_{\max} = \frac{w_0 L^2}{8} = \frac{12\,000 \times 6^2}{8}
= 54\,000\,\text{N}\cdot\text{m} = 54\,\text{kN}\cdot\text{m}.
\]
With an allowable bending stress $\sigma_{\text{allow}} =
\sigma_y/1.67 = 250/1.67 \approx 150\,\text{MPa}$ for ASD on mild
steel, the strength requirement is
\[
S \geq \frac{M_{\max}}{\sigma_{\text{allow}}} = \frac{54\,000}{150
\times 10^6} = 3.6 \times 10^{-4}\,\si{\meter\cubed} = 360\,\text{cm}^3.
\]
The deflection requirement, with $\delta \leq L/360 = 6000/360 =
16.7\,\si{\milli\meter}$, is
\[
\frac{5 w_0 L^4}{384 EI} \leq \frac{L}{360}
\;\Rightarrow\;
I \geq \frac{360 \times 5 w_0 L^3}{384 E}
= \frac{1800 \times 12\,000 \times 6^3}{384 \times 200 \times 10^9}
= 6.08 \times 10^{-5}\,\si{\meter\tothe{4}} = 6080\,\text{cm}^4.
\]
A $W12\times 65$ ($S = 1440\,\text{cm}^3$, $I = 22\,200\,\text{cm}^4$)
satisfies both with wide margin; a lighter $W16\times 40$
($S=1070\,\text{cm}^3$, $I=21\,800\,\text{cm}^4$) also satisfies both
and is the economical choice, since both its strength and stiffness
exceed the requirements while it weighs less. The deeper, lighter
section wins, which is the recurring lesson that depth buys stiffness
cheaply.

\paragraph{Example 2: a pinned steel column.}
A pinned-pinned column of unbraced length $L = 4\,\si{\meter}$ carries
an axial service load of $300\,\text{kN}$. Using a $W10\times 49$
($A = 92.9\,\text{cm}^2$, weak-axis $r_y = 65\,\si{\milli\meter}$,
$I_y = A r_y^2 = 92.9\times 10^{-4} \times 0.065^2 = 3.93\times
10^{-5}\,\si{\meter\tothe{4}}$), check the Euler buckling capacity
about the governing weak axis.

The slenderness about the weak axis is
\[
\lambda = \frac{KL}{r_y} = \frac{1.0 \times 4000}{65} = 61.5,
\]
below the transition $\lambda_c \approx 89$, so the column is in the
intermediate range where neither pure Euler buckling nor pure yielding
governs. The elastic Euler load is
\[
P_{\text{cr}} = \frac{\pi^2 E I_y}{(KL)^2}
= \frac{\pi^2 \times 200 \times 10^9 \times 3.93 \times 10^{-5}}
{4^2}
= 4.85 \times 10^6\,\text{N} = 4850\,\text{kN},
\]
and the squash (yield) load is $P_y = \sigma_y A = 250\times 10^6
\times 92.9 \times 10^{-4} = 2320\,\text{kN}$. The intermediate-column
capacity lies below both; a code column curve in this slenderness
range gives roughly $0.8 P_y \approx 1860\,\text{kN}$, comfortably
above the $300\,\text{kN}$ demand. The column is governed by yielding
with an imperfection allowance, not by elastic buckling, exactly as the
slenderness $\lambda < \lambda_c$ predicted.

\paragraph{Example 3: combined bending and torsion in a shaft.}
A solid circular shaft of diameter $d = 40\,\si{\milli\meter}$ carries
a bending moment $M = 600\,\text{N}\cdot\text{m}$ and a torque $T =
500\,\text{N}\cdot\text{m}$. Find the maximum principal stress and the
maximum shear stress at the surface.

For a solid circle, $I = \pi d^4/64$ and $J = \pi d^4/32 = 2I$, so the
bending and torsional section moduli are $S = I/(d/2) = \pi d^3/32$ and
$J/(d/2) = \pi d^3/16 = 2S$. With $d = 0.040\,\si{\meter}$,
\[
\pi d^3/32 = \frac{\pi \times 0.040^3}{32} = 6.28 \times 10^{-6}
\,\si{\meter\cubed}.
\]
The bending and torsional surface stresses are
\[
\sigma = \frac{M}{\pi d^3/32} = \frac{600}{6.28\times 10^{-6}}
= 95.5\,\text{MPa}, \qquad
\tau = \frac{T}{\pi d^3/16} = \frac{500}{2 \times 6.28\times 10^{-6}}
= 39.8\,\text{MPa}.
\]
At the surface the stress state is $\sigma_{xx} = 95.5\,\text{MPa}$,
$\sigma_{yy} = 0$, $\sigma_{xy} = 39.8\,\text{MPa}$. The principal
stresses from the Mohr's circle of chapter 3 are
\[
\sigma_{1,2} = \frac{\sigma_{xx}}{2} \pm \sqrt{\left(\frac{\sigma_{xx}}
{2}\right)^2 + \tau^2}
= 47.7 \pm \sqrt{47.7^2 + 39.8^2}
= 47.7 \pm 62.1\,\text{MPa},
\]
giving $\sigma_1 = 109.8\,\text{MPa}$ (the maximum principal stress)
and the maximum in-plane shear stress $\tau_{\max} = 62.1\,\text{MPa}$.
A von Mises check, $\sigma_{vM} = \sqrt{\sigma_{xx}^2 + 3\tau^2} =
\sqrt{95.5^2 + 3\times 39.8^2} = 116\,\text{MPa}$, would be compared to
the yield strength to size the shaft.

\paragraph{Example 4: deflection of an overhanging beam by superposition.}
A simply supported beam of span $L$ carries a central point load $P$
and also a uniform line load $w_0$ over the full span. Find the midspan
deflection.

Superposition adds the two single-load results directly, since the
beam is linear elastic:
\[
\delta_{\text{mid}} = \frac{PL^3}{48EI} + \frac{5 w_0 L^4}{384 EI}.
\]
With $P = 20\,\text{kN}$, $w_0 = 5\,\text{kN}/\text{m}$, $L =
8\,\si{\meter}$, $E = 200\,\text{GPa}$, and a $W12\times 65$
($I = 22\,200\,\text{cm}^4 = 2.22\times 10^{-4}\,\si{\meter\tothe{4}}$):
\[
\delta_{\text{mid}} = \frac{20\,000 \times 8^3}{48 \times 200\times
10^9 \times 2.22\times 10^{-4}} + \frac{5 \times 5000 \times 8^4}
{384 \times 200\times 10^9 \times 2.22\times 10^{-4}}.
\]
The two terms are $4.81\,\si{\milli\meter}$ and $6.01\,\si{\milli
\meter}$, summing to $\delta_{\text{mid}} = 10.8\,\si{\milli\meter}$.
This matches, to within the discretisation error, the moment computed
by \path{code/shear_moment.py} for the same loading, whose midspan
moment $54 + 40 = 94\,\text{kN}\cdot\text{m}$ feeds the bending-stress
check in parallel.

\paragraph{Example 5: lateral-torsional buckling sets the unbraced length.}
A simply supported steel beam carries a moment that would be safe in
bending if the compression flange were braced. With the flange
unbraced over the full span, a code lateral-torsional buckling check
reduces the allowable moment. Suppose the fully braced (plastic)
moment is $M_p = 200\,\text{kN}\cdot\text{m}$ and the code gives a
limiting unbraced length $L_p = 2.1\,\si{\meter}$ below which the full
$M_p$ is available, and $L_r = 6.3\,\si{\meter}$ above which elastic
lateral-torsional buckling governs. For an unbraced length $L_b =
4.2\,\si{\meter}$ between these, the AISC inelastic-buckling
interpolation gives
\[
M_n = M_p\left[1 - 0.3\,\frac{L_b - L_p}{L_r - L_p}\right]
= 200\left[1 - 0.3\,\frac{4.2 - 2.1}{6.3 - 2.1}\right]
= 200(1 - 0.15) = 170\,\text{kN}\cdot\text{m}.
\]
Bracing the compression flange at midspan, halving $L_b$ to
$2.1\,\si{\meter} = L_p$, restores the full $M_p = 200\,\text{kN}\cdot
\text{m}$: a single brace recovers eighteen per cent of capacity at
negligible cost. This is the calculation behind the rule that floor
beams must be braced at intervals, and the reason an unbraced beam
during erection is weaker than the same beam in service.

\section{An end-to-end floor-beam design check}\pathtag{standard}

The methods of the chapter are exercised one at a time in the worked
examples. A real design check runs them in sequence on a single member,
carrying the same numbers from the load path through to the verdict, and
the sequence is itself the discipline: a member that passes one check
and fails another is not designed until every check has been made. The
case study below sizes an interior steel floor beam in an office
building, from the floor pressure down to the connection note, and
records at each stage which check governs.

\subsection{The brief and the load path}

A typical office floor bay frames as a grid of joists spanning onto
interior beams, which in turn span onto girders. We design one interior
beam. The beam spans $L = 8\,\si{\meter}$ between girder centrelines and
sits at a spacing such that its tributary width is $w_t = 3\,\si{\meter}$
(half the joist span on each side). The floor carries a superimposed
dead load (slab, finishes, services, partitions) of $g = 4.0\,\text{kN}/
\text{m}^2$ and an office live load of $q = 3.0\,\text{kN}/\text{m}^2$,
the latter a standard occupancy value. Steel is grade $S355$
($\sigma_y = 355\,\text{MPa}$, $E = 200\,\text{GPa}$).

The first step is the load path of figure
\ref{fig:vol03ch08:floor-beam-loadpath}: floor pressure collected over
the tributary width becomes a line load on the beam. The service line
loads are
\[
g_0 = g\,w_t = 4.0 \times 3 = 12\,\text{kN}/\text{m}, \qquad
q_0 = q\,w_t = 3.0 \times 3 = 9\,\text{kN}/\text{m}.
\]
The beam self-weight, not yet known, is small against these and is added
once a trial section is on the table.

% Figure: tributary-area load path for the case-study floor beam.
% A row of joists at spacing s delivers a uniform line load to the
% interior beam over span L; the beam delivers half its load to each
% of two girders. The diagram shows the tributary width and the
% resulting line load.
\begin{figure}[htbp]
\centering
\begin{tikzpicture}[scale=0.9]
% plan view of a floor bay
% the interior beam (heavy) runs left-right
\draw[engblack, very thick] (0,0) -- (8,0);
\node[engblack, anchor=north, font=\scriptsize] at (4,-0.15) {interior beam, span $L$};
% girders at the two ends, running into the page (drawn vertical)
\draw[engblue, line width=3pt] (0,-1.6) -- (0,1.6);
\draw[engblue, line width=3pt] (8,-1.6) -- (8,1.6);
\node[engblue, anchor=east, font=\scriptsize] at (-0.1,1.4) {girder};
\node[engblue, anchor=west, font=\scriptsize] at (8.1,1.4) {girder};
% joists running parallel to the girders, spaced along the beam
\foreach \x in {1,2,3,4,5,6,7}{
  \draw[enggray, thick] (\x,-1.3) -- (\x,1.3);
}
\node[enggray, anchor=south, font=\scriptsize] at (3.5,1.3) {joists at spacing $s$};
% tributary band of the interior beam (half a joist span each side)
\fill[engamber!18] (0,-0.8) rectangle (8,0.8);
\draw[<->, engamber] (8.6,-0.8) -- (8.6,0.8);
\node[engamber, anchor=west, font=\scriptsize] at (8.7,0) {tributary width $w_t$};
% resultant line load arrows on the beam
\foreach \x in {0.5,1.5,2.5,3.5,4.5,5.5,6.5,7.5}{
  \draw[engred, ->] (\x,1.1) -- (\x,0.12);
}
\node[engred, anchor=south, font=\scriptsize] at (4,1.15) {line load $w_0 = q\,w_t$};
\end{tikzpicture}
\caption[Tributary load path for a floor beam]{Plan of one floor bay.
The floor pressure $q$ is collected by the joists and delivered to the
interior beam over a tributary width $w_t$ equal to half the joist span
on each side of the beam. The beam therefore carries a uniform line
load $w_0 = q\,w_t$, and passes half of its total load to each of the
two girders at its ends. The first step of every floor-beam design is
this load path: pressure to tributary width to line load to end
reactions, before any stress or deflection is computed.}
\label{fig:vol03ch08:floor-beam-loadpath}
\end{figure}


\subsection{Factored load and the demand}

Limit-state design factors the loads up. Using the load combination
$1.35 g + 1.5 q$ (the Eurocode persistent-design combination, close to
the AISC $1.2D + 1.6L$),
\[
w_u = 1.35\,g_0 + 1.5\,q_0 = 1.35 \times 12 + 1.5 \times 9
= 16.2 + 13.5 = 29.7\,\text{kN}/\text{m}.
\]
The simply supported beam carries the demand moment and shear
\[
M_u = \frac{w_u L^2}{8} = \frac{29.7 \times 8^2}{8}
= 237.6\,\text{kN}\cdot\text{m}, \qquad
V_u = \frac{w_u L}{2} = \frac{29.7 \times 8}{2} = 118.8\,\text{kN}.
\]
These are the numbers every subsequent check is measured against.

\subsection{Section selection by required modulus}

The flexural strength requirement, with a resistance factor $\phi_b =
0.9$ and the plastic section modulus $Z$ governing for a compact
section, is $\phi_b \sigma_y Z \geq M_u$, so
\[
Z_{\text{req}} = \frac{M_u}{\phi_b \sigma_y}
= \frac{237.6 \times 10^3}{0.9 \times 355 \times 10^6}
= 7.44 \times 10^{-4}\,\si{\meter\cubed} = 744\,\text{cm}^3.
\]
Scanning the section table, a $W16\times 40$ ($S = 1070\,\text{cm}^3$,
$I = 21\,800\,\text{cm}^4$, depth $407\,\si{\milli\meter}$) carries a
plastic modulus of roughly $Z \approx 1.12\,S \approx 1200\,\text{cm}^3$
for a rolled I-section, comfortably above the $744\,\text{cm}^3$ demand.
The strength check passes at the trial section with a utilisation of
about $0.62$. We carry the $W16\times 40$ forward and add its
self-weight, about $0.58\,\text{kN}/\text{m}$, which lifts $w_u$ by
roughly $0.8\,\text{kN}/\text{m}$, a two-and-a-half per cent increase
that the margin absorbs.

\subsection{Shear check}

The web carries the shear. For a $W16\times 40$ the web area is
approximately $A_w = d\,t_w \approx 0.407 \times 0.0078 = 3.17 \times
10^{-3}\,\si{\meter\squared}$. The shear capacity is $\phi_v\,0.6
\sigma_y A_w$ with $\phi_v = 0.9$:
\[
\phi_v V_n = 0.9 \times 0.6 \times 355 \times 10^6 \times 3.17 \times
10^{-3} = 608\,\text{kN},
\]
far above the $118.8\,\text{kN}$ demand. Shear rarely governs a
rolled-section floor beam of normal span, and it does not here; the
utilisation is below $0.2$.

\subsection{Deflection serviceability}

Serviceability is checked under unfactored loads. The live-load
deflection, the quantity occupants perceive and the one that cracks
finishes, is
\[
\delta_{q} = \frac{5\,q_0 L^4}{384\,E I}
= \frac{5 \times 9000 \times 8^4}{384 \times 200 \times 10^9 \times
2.18 \times 10^{-4}}
= 11.0\,\si{\milli\meter}.
\]
The limit for a floor supporting non-brittle finishes is $L/360 =
8000/360 = 22.2\,\si{\milli\meter}$, so the beam passes with a
utilisation of $0.50$. The total-load deflection, $\delta_{g+q}$ under
$g_0 + q_0 = 21\,\text{kN}/\text{m}$, is $25.7\,\si{\milli\meter}$
against a looser $L/240 = 33.3\,\si{\milli\meter}$ limit, which also
passes. Deflection is the closest check so far, which is the usual
finding for a beam of this span; a shallower section would have been
governed by it.

\subsection{Lateral-torsional buckling}

The flexural capacity assumed the compression flange braced. The floor
slab cast onto the top flange braces it continuously in the finished
condition, so $L_b \approx 0$ and the full plateau moment is available;
the strength check above stands. The check that matters is the
construction stage, before the slab is poured, when the beam may be
braced only at the points where joists frame in. Suppose joists land at
the third points, giving an unbraced length $L_b = 8/3 = 2.67\,
\si{\meter}$. For a $W16\times 40$ the plateau limit is about $L_p =
1.7\,\si{\meter}$ and the elastic limit about $L_r = 5.2\,
\si{\meter}$, so $L_b$ falls in the inelastic range. The interpolation
of figure \ref{fig:vol03ch08:ltb-capacity},
\[
M_n = M_p\left[1 - 0.3\,\frac{L_b - L_p}{L_r - L_p}\right]
= M_p\left[1 - 0.3\,\frac{2.67 - 1.7}{5.2 - 1.7}\right]
= 0.917\,M_p,
\]
keeps about ninety-two per cent of capacity during erection, which the
strength margin still covers. The note on the drawing is that the
joists must be in place before the construction live load is applied,
and that no point load may be hung from the unbraced midspan during
erection.

% Figure: nominal flexural capacity M_n versus unbraced length L_b.
% Plateau at M_p up to L_p, linear inelastic LTB segment from L_p to L_r,
% elastic LTB hyperbola beyond L_r. Numbers match the case study:
% M_p = 175, M_r = 0.7 M_p = 122.5, L_p = 2.0, L_r = 6.0 (schematic units).
\begin{figure}[htbp]
\centering
\begin{tikzpicture}
\begin{axis}[
  width=10cm, height=6.2cm,
  xlabel={unbraced length $L_b$ (m)},
  ylabel={nominal moment $M_n$ (kN$\cdot$m)},
  xmin=0, xmax=9, ymin=0, ymax=210,
  xtick={0,2,4,6,8},
  ytick={0,50,100,150,200},
  axis lines=left,
  every axis plot/.append style={very thick},
  legend style={font=\scriptsize, at={(0.97,0.95)}, anchor=north east},
  tick label style={font=\scriptsize},
  label style={font=\scriptsize},
]
% plateau: full plastic moment up to L_p = 2
\addplot[engblue] coordinates {(0,175) (2,175)};
\addlegendentry{plateau $M_p$}
% inelastic LTB: linear from (2,175) to (6, 122.5)
\addplot[engamber] coordinates {(2,175) (6,122.5)};
\addlegendentry{inelastic LTB}
% elastic LTB: hyperbola-like decay beyond L_r=6, M ~ C / L_b^2 matched at (6,122.5)
\addplot[engred, domain=6:9, samples=40] {122.5*36/(x*x)};
\addlegendentry{elastic LTB}
% markers for L_p and L_r
\addplot[only marks, mark=*, engblack, mark size=1.6pt] coordinates {(2,175) (6,122.5)};
\node[anchor=south, font=\scriptsize] at (axis cs:2,178) {$L_p$};
\node[anchor=south west, font=\scriptsize] at (axis cs:6,124) {$L_r$};
\end{axis}
\end{tikzpicture}
\caption[Flexural capacity versus unbraced length]{The nominal
flexural capacity of a steel beam falls as the unbraced length of its
compression flange grows. Below $L_p$ the section reaches its full
plastic moment $M_p$ on the plateau. Between $L_p$ and $L_r$ the
capacity drops linearly as inelastic lateral-torsional buckling begins
to interact with yielding. Beyond $L_r$ the capacity follows the
elastic lateral-torsional buckling curve, decaying roughly with the
square of the unbraced length. The design lesson is direct: a brace
that shortens $L_b$ moves the design left along the curve and recovers
capacity, which is why floor beams are braced at intervals and why the
same beam is weaker during erection than in service.}
\label{fig:vol03ch08:ltb-capacity}
\end{figure}


\subsection{Connection note and verdict}

The beam delivers its end reaction $V_u = 118.8\,\text{kN}$ (factored)
to the girder through a simple shear connection: a bolted double-angle
or a shear tab, designed for the reaction with no moment, since the
analysis modelled the end as a pin. The connection check is its own
calculation, against bolt shear, bolt bearing on the beam web, and
block shear of the connected web, each verified separately; a pair of
M20 bolts in grade 8.8 carries roughly $94\,\text{kN}$ each in single
shear, so three bolts clear the reaction with margin. The connection is
detailed and checked as a subassembly, exactly the joint check whose
omission caused the Hyatt Regency collapse discussed below.

The verdict: a $W16\times 40$ in $S355$ steel carries the office floor
bay with the governing check being live-load deflection at a
utilisation of $0.50$, flexure at $0.62$, and every other mode well
inside its limit. The member is governed by serviceability, not
strength, and the design is signed off on that basis. The full table of
utilisations is the deliverable that an independent checker reviews; a
single governing number hides the six other checks that had to pass for
it to be the governing one.

\begin{principle}[A design check is the full sequence, not one number]
A structural member is designed only when every applicable limit state
has been verified against the same set of loads: load path, factored
demand, flexural strength, shear, deflection serviceability, stability
during construction and service, and the end connection. The governing
check is whichever has the highest utilisation, and it is meaningful
only because the others were made and passed. Reporting strength alone,
or stopping at the first comfortable margin, is the habit that lets a
serviceability failure or an unbraced-erection buckling reach the field.
\end{principle}

\section{Further worked examples}\pathtag{standard}

The end-to-end case study runs the checks in series on one member. The
three examples here isolate methods that the case study did not call on:
an indeterminate beam solved by the force method, a continuous beam by
moment distribution, and a beam-column under combined axial and bending
load checked against the interaction envelope.

\paragraph{Example 6: a propped cantilever by the force method.}
A propped cantilever of span $L = 6\,\si{\meter}$ carries a uniform load
$w_0 = 20\,\text{kN}/\text{m}$. The structure is fixed at the left and
simply propped at the right, once statically indeterminate. Find the
prop reaction, the fixed-end moment, and the maximum sagging moment in
the span.

Choose the prop reaction $R$ as the redundant. Releasing the prop
leaves a cantilever fixed at the left and free at the right, the primary
structure of figure \ref{fig:vol03ch08:propped-cantilever}. Two
deflections at the released point matter: the downward tip deflection
under the distributed load, $\delta_w = w_0 L^4/(8EI)$, and the upward
tip deflection under the unknown prop reaction acting alone, $\delta_R =
R L^3/(3EI)$. Compatibility requires the propped end not to move, so
$\delta_w = \delta_R$:
\[
\frac{w_0 L^4}{8 EI} = \frac{R L^3}{3 EI}
\;\Rightarrow\;
R = \frac{3 w_0 L}{8} = \frac{3 \times 20 \times 6}{8} = 45\,\text{kN}.
\]
The fixed-end moment follows from statics on the released cantilever
with the prop reaction restored, taking moments about the fixed end:
\[
M_{\text{fix}} = \frac{w_0 L^2}{2} - R L
= \frac{20 \times 6^2}{2} - 45 \times 6 = 360 - 270 = 90\,\text{kN}
\cdot\text{m},
\]
a hogging moment of magnitude $w_0 L^2/8 = 90\,\text{kN}\cdot\text{m}$,
the closed-form result for this case. The shear is zero where the
sagging moment peaks; measuring $x$ from the propped end, $V(x) = R -
w_0 x = 0$ at $x = R/w_0 = 45/20 = 2.25\,\si{\meter}$, and the sagging
moment there is
\[
M_{\text{span}} = R x - \frac{w_0 x^2}{2}
= 45 \times 2.25 - \frac{20 \times 2.25^2}{2}
= 101.25 - 50.6 = 50.6\,\text{kN}\cdot\text{m}.
\]
The fixed end at $90\,\text{kN}\cdot\text{m}$ is the critical section,
half again the span peak. The redistribution the prop forces, hogging
at the wall and sagging in the span, is what makes the indeterminate
member carry more than a cantilever of the same span and stiffer than a
simple beam.

\paragraph{Example 7: a two-span continuous beam by moment distribution.}
A continuous beam runs over three supports, with two equal spans of
$L = 6\,\si{\meter}$ each, carrying a uniform load $w_0 = 15\,\text{kN}/
\text{m}$ on both spans. Find the moment over the central support.

Moment distribution (the Hardy Cross method, treated in the mastery box
below) starts from the fixed-end moments, the moments that would arise
if every joint were clamped. For a uniformly loaded span clamped at both
ends the fixed-end moments are $\pm w_0 L^2/12$:
\[
M_F = \frac{w_0 L^2}{12} = \frac{15 \times 6^2}{12} = 45\,\text{kN}
\cdot\text{m}.
\]
At the central support two spans meet, each contributing a fixed-end
moment of $45\,\text{kN}\cdot\text{m}$ but of opposite sense, so before
release they nearly balance; the unbalanced moment at the central joint
comes from the difference in carry-over from the far (simply supported)
ends. The two end supports are simple, so their fixed-end moments
release fully and carry over one half to the central support. Working
the single distribution cycle for the symmetric two-span beam gives the
standard result for equal spans and equal load,
\[
M_{\text{central}} = -\frac{w_0 L^2}{8} = -\frac{15 \times 6^2}{8}
= -67.5\,\text{kN}\cdot\text{m},
\]
a hogging moment over the central support larger in magnitude than the
$w_0 L^2/8 = 67.5\,\text{kN}\cdot\text{m}$ a single simple span would
carry at midspan, which is why the support region of a continuous beam,
not its midspan, usually governs. The midspan sagging moment in each
span drops to $9 w_0 L^2/128 \approx 38\,\text{kN}\cdot\text{m}$,
roughly nine-sixteenths of the simple-span value, the saving continuity
buys at the cost of a heavier support moment.

\paragraph{Example 8: a beam-column under combined axial and bending.}
A column in a braced frame carries a factored axial compression $P_r =
600\,\text{kN}$ and, because it also supports a beam framing in with an
eccentric reaction, a factored moment $M_r = 80\,\text{kN}\cdot\text{m}$
about its strong axis. Its axial capacity for the governing slenderness
is $P_c = 1700\,\text{kN}$ and its flexural capacity is $M_c =
250\,\text{kN}\cdot\text{m}$. Check the member.

The axial ratio is $P_r/P_c = 600/1700 = 0.35$, above the $0.2$ kink, so
the steeper branch of the AISC interaction equation governs:
\[
\frac{P_r}{P_c} + \frac{8}{9}\,\frac{M_r}{M_c}
= 0.35 + \frac{8}{9}\times\frac{80}{250}
= 0.35 + \frac{8}{9}\times 0.32 = 0.35 + 0.284 = 0.64 \leq 1.0.
\]
The member passes at a combined utilisation of $0.64$, plotted as the
demand point inside the envelope of figure
\ref{fig:vol03ch08:beam-column-interaction}. Neither load alone is near
its capacity ($0.35$ axial, $0.32$ flexural), yet the interaction
consumes nearly two-thirds of the member, the point of the combined
check: a member acceptable under each load separately can be overloaded
by their combination, and the envelope is what catches it. Were the
axial ratio below $0.2$ the shallower branch $P_r/(2P_c) + M_r/M_c \leq
1$ would apply, weighting bending more heavily, because in the
bending-dominated regime the moment is the load that drives the member
to its limit.

% Figure: propped cantilever under uniform load. Fixed end at left,
% simple prop at right. Shows the released primary structure (cantilever)
% and the redundant prop reaction R, plus the resulting bending-moment
% diagram with the hogging fixed-end moment wL^2/8 and the sagging span peak.
\begin{figure}[htbp]
\centering
\begin{tikzpicture}[scale=1.0]
% structure sketch
\begin{scope}[yshift=3.0cm]
  % wall (fixed end) at left
  \draw[engblack, line width=2pt] (0,-0.6) -- (0,0.6);
  \foreach \y in {-0.6,-0.35,-0.1,0.15,0.4}{ \draw[enggray] (0,\y) -- (-0.25,\y-0.2);}
  % beam
  \draw[engblack, very thick] (0,0) -- (6,0);
  % uniform load
  \foreach \x in {0.3,0.9,1.5,2.1,2.7,3.3,3.9,4.5,5.1,5.7}{
    \draw[engred, ->] (\x,0.85) -- (\x,0.1);
  }
  \draw[engred, thick] (0,0.85) -- (6,0.85);
  \node[engred, anchor=south, font=\scriptsize] at (3,0.9) {$w_0$};
  % prop (roller) at right
  \fill[engblue] (6,0) circle (0.11);
  \draw[engblue] (5.6,-0.42) -- (6.4,-0.42);
  \node[engblue, anchor=north, font=\scriptsize] at (6,-0.45) {prop $R$};
  \node[engblack, anchor=north east, font=\scriptsize] at (-0.05,-0.2) {fixed};
  \draw[<->, enggray] (0,-0.85) -- (6,-0.85) node[midway, below, font=\scriptsize]{span $L$};
\end{scope}
% bending-moment diagram
\begin{scope}[yshift=0.2cm]
  \draw[enggray, ->] (-0.3,0) -- (6.5,0) node[anchor=west, font=\scriptsize]{$x$};
  \node[anchor=east, font=\scriptsize] at (-0.3,0.5) {$M(x)$};
  % hogging fixed-end moment at left (drawn below axis), magnitude wL^2/8
  % sagging hump in span (above axis), peak near 0.4L
  % build a smooth-ish piecewise curve with samples
  \draw[engred, very thick]
    (0,-1.1)
    .. controls (1.0,-0.2) and (1.6,0.45) .. (2.4,0.62)
    .. controls (3.4,0.82) and (4.6,0.55) .. (6,0);
  \fill[engred!12]
    (0,-1.1)
    .. controls (1.0,-0.2) and (1.6,0.45) .. (2.4,0.62)
    .. controls (3.4,0.82) and (4.6,0.55) .. (6,0) -- (0,0) -- cycle;
  \node[engred, anchor=north east, font=\scriptsize] at (0.1,-1.1) {$M_{\text{fix}}=w_0L^2/8$ (hog)};
  \node[engred, anchor=south, font=\scriptsize] at (2.7,0.82) {$M_{\text{span}}\approx 9w_0L^2/128$ (sag)};
\end{scope}
\end{tikzpicture}
\caption[Propped cantilever, moment diagram]{A propped cantilever
under uniform load: fixed at the left, simply propped at the right. The
structure is once indeterminate; releasing the prop leaves a
cantilever, and the prop reaction is the redundant fixed by the
compatibility condition that the propped end does not deflect. The
moment diagram carries a hogging peak $w_0L^2/8$ at the fixed end and a
smaller sagging peak of about $9w_0L^2/128$ in the span, located at
$x=5L/8$ from the fixed end. The fixed end is the critical section for
this member.}
\label{fig:vol03ch08:propped-cantilever}
\end{figure}


% Figure: beam-column interaction diagram. Axial ratio on the vertical
% axis, moment ratio on the horizontal. Shows the AISC bilinear envelope
% (steeper slope for P/Pc >= 0.2, shallower below) and a sample demand
% point inside the safe region.
\begin{figure}[htbp]
\centering
\begin{tikzpicture}
\begin{axis}[
  width=8.4cm, height=7.2cm,
  xlabel={moment ratio $M_r/M_c$},
  ylabel={axial ratio $P_r/P_c$},
  xmin=0, xmax=1.15, ymin=0, ymax=1.15,
  xtick={0,0.2,0.4,0.6,0.8,1.0},
  ytick={0,0.2,0.4,0.6,0.8,1.0},
  axis lines=left,
  tick label style={font=\scriptsize},
  label style={font=\scriptsize},
  legend style={font=\scriptsize, at={(0.97,0.97)}, anchor=north east},
]
% AISC H1-1 bilinear envelope:
% for P/Pc >= 0.2: P/Pc + (8/9)(M/Mc) = 1
% for P/Pc <  0.2: P/Pc/2 + (M/Mc) = 1
% upper segment: from (M=0,P=1) to the kink at P=0.2
% at P=0.2: M/Mc = (1-0.2)*9/8 = 0.9
\addplot[engblue, very thick] coordinates {(0,1.0) (0.9,0.2)};
% lower segment: from kink (0.9,0.2) to (M=1, P=0)
% check: P/Pc/2 + M/Mc = 1 -> at P=0.2: 0.1 + M =1 -> M=0.9 OK; at P=0 -> M=1
\addplot[engblue, very thick] coordinates {(0.9,0.2) (1.0,0.0)};
\addlegendentry{AISC envelope}
% safe-region shading
\addplot[engblue!10, draw=none] coordinates {(0,0) (0,1.0) (0.9,0.2) (1.0,0.0)} \closedcycle;
% sample demand point inside
\addplot[only marks, mark=*, engred, mark size=2pt] coordinates {(0.45,0.35)};
\node[engred, anchor=west, font=\scriptsize] at (axis cs:0.47,0.37) {demand};
% kink marker
\addplot[only marks, mark=*, engblack, mark size=1.2pt] coordinates {(0.9,0.2)};
\node[engblack, anchor=south west, font=\scriptsize] at (axis cs:0.82,0.22) {kink at $0.2$};
\end{axis}
\end{tikzpicture}
\caption[Beam-column interaction diagram]{The interaction envelope for
a member carrying combined axial compression and bending. The vertical
axis is the required axial force as a fraction of the axial capacity,
the horizontal axis the required moment as a fraction of the moment
capacity. The bilinear envelope steepens above an axial ratio of
$0.2$, where the axial term dominates, and flattens below it, where
bending dominates. Any demand point inside the envelope passes the
combined check; a point outside fails even though each load taken alone
might be acceptable. The kink reflects the design judgement that the
two failure modes interact differently depending on which load is the
larger.}
\label{fig:vol03ch08:beam-column-interaction}
\end{figure}


\begin{mastery}
\textbf{Moment distribution: the Hardy Cross method.} Before the
stiffness method ran on every laptop, the analysis of continuous beams
and frames was carried out by hand with the iterative scheme Hardy Cross
published in 1932, and it remains the fastest hand check for a
continuous beam and the clearest picture of how moment redistributes
through a rigid-jointed structure. The method works with three member
properties. The \engterm{stiffness} of a prismatic member fixed at the
far end is $k = 4EI/L$; if the far end is pinned it is $k = 3EI/L$. The
\engterm{distribution factor} at a joint apportions an unbalanced moment
among the members meeting there in proportion to their stiffnesses,
$\text{DF}_i = k_i / \sum_j k_j$, and the factors at a joint sum to one.
The \engterm{carry-over factor} sends half of any moment applied at one
end of a prismatic member to its fixed far end, carry-over $= 1/2$.

The procedure is a relaxation. Clamp every joint and write the
fixed-end moments for the applied loads (for a uniform load on a span,
$\pm w_0 L^2/12$; for a central point load, $\pm PL/8$). Each joint
then carries an unbalanced moment, the algebraic sum of the fixed-end
moments of the members meeting there. Release one joint at a time: the
unbalanced moment is distributed to the members in proportion to their
distribution factors, with the sign that restores balance, and half of
each distributed moment is carried over to the member's far end. The
carry-overs unbalance the neighbouring joints, which are released in
turn, and the process repeats. Because each cycle multiplies the
residual unbalanced moments by a factor well below one, the residuals
shrink geometrically and the moments converge to the exact
elastic answer in a handful of cycles. Summing the fixed-end moment, the
distributed moments, and the carried-over moments at each member end
gives the final moment there.

The two-span continuous beam of example 7 falls out of a single
distribution cycle by symmetry; an asymmetric three-span beam or a
no-sway frame needs three or four cycles to converge to three
significant figures. Sidesway frames need an extra layer, in which the
sway moments are found by releasing the frame to translate, then
superposed in the proportion that satisfies horizontal equilibrium. The
method is no longer the production tool, the stiffness method of the
mastery box in section 8.5 is, yet an engineer who can run a moment
distribution on the back of a drawing can check a software model's
support moments in two minutes and catch the gross input error that a
column of finite-element output would hide.
\end{mastery}

\section{Failure: column buckling, beam fatigue, joint detail
failures}\pathtag{core}

Structural failures in beams, columns, and frames cluster into
three families: instability of compression members, fatigue of
flexural members, and detail-level failures at connections. The
chapter's failure section names the canonical mechanisms.

\subsection{Column buckling}

A column whose slenderness ratio exceeds the transition value
$\lambda_c$ fails by Euler buckling at a load below the material's
compressive strength. The failure is sudden: at $P < P_{\text{cr}}$
the column is straight and in equilibrium; at $P > P_{\text{cr}}$
the column deflects laterally, and the deflection grows with no
external warning until either material yielding limits it or the
column collapses.

% Figure: post-buckling pitchfork bifurcation. Load P/Pcr vs midspan
% lateral deflection. Straight branch up to Pcr, then two symmetric
% deflected branches rising slightly (supercritical pitchfork).
\begin{figure}[htbp]
\centering
\begin{tikzpicture}
\begin{axis}[
  width=11cm, height=7cm,
  xlabel={midspan lateral deflection $\delta/L$},
  ylabel={axial load $P/P_{\mathrm{cr}}$},
  xmin=-0.4, xmax=0.4, ymin=0, ymax=1.6,
  axis lines=middle,
  tick label style={font=\scriptsize}, label style={font=\small},
  ytick={1.0}, yticklabels={$1$},
  xtick={0},
]
% trivial (straight) branch, stable below Pcr, unstable above
\addplot[engblue, very thick, domain=0:1.0, samples=2] coordinates {(0,0) (0,1.0)};
\addplot[engblue, very thick, dashed, domain=1.0:1.6, samples=2] coordinates {(0,1.0) (0,1.6)};
% post-buckling branches: P/Pcr = 1 + (pi^2/8)(delta/L)^2 approx
\addplot[engred, very thick, domain=0:0.38, samples=60]
  {1 + (pi^2/8)*x^2};
\addplot[engred, very thick, domain=-0.38:0, samples=60]
  {1 + (pi^2/8)*x^2};
% bifurcation point
\addplot[engblack, only marks, mark=*, mark options={fill=engamber}]
  coordinates {(0,1.0)};
\node[engamber, font=\scriptsize, anchor=south west] at (axis cs:0.02,1.02)
  {bifurcation $P=P_{\mathrm{cr}}$};
\node[engblue, font=\scriptsize, anchor=east] at (axis cs:-0.01,0.5) {straight (stable)};
\node[engblue, font=\scriptsize, anchor=west] at (axis cs:0.01,1.35) {straight (unstable)};
\node[engred, font=\scriptsize, anchor=south west] at (axis cs:0.2,1.18) {buckled branch};
\end{axis}
\end{tikzpicture}
\caption[Supercritical pitchfork at the buckling load]{Load versus
lateral deflection for a pinned-pinned column, showing the
supercritical pitchfork bifurcation at the Euler load. Below
$P_{\mathrm{cr}}$ the only equilibrium is the straight column (blue,
stable). At $P=P_{\mathrm{cr}}$ the straight state loses stability
(blue dashed) and two symmetric deflected branches appear (red),
along which the column carries slightly more than $P_{\mathrm{cr}}$
as the deflection grows. Linear Euler theory predicts the
bifurcation load but not the post-buckling path; the gentle rise of
the red branches is a result of the nonlinear large-deflection
(elastica) theory. Real columns, with imperfections, follow a
rounded curve that approaches the bifurcation diagram from one side
and never reaches the ideal straight branch.}
\label{fig:vol03ch08:postbuckling}
\end{figure}


The mathematics behind the suddenness is the bifurcation diagram of
figure \ref{fig:vol03ch08:postbuckling}: at $P_{\text{cr}}$ the
straight equilibrium exchanges stability with a pair of deflected
branches in a supercritical pitchfork. Linear theory locates the
bifurcation but says nothing about the path beyond it; the
large-deflection (elastica) theory shows the deflected branches rising
only gently above $P_{\text{cr}}$, so a real column near the critical
load has almost no reserve. The standard examples:

\begin{itemize}[leftmargin=*]
  \item Long compression chords in trusses; a slender wood or steel
    chord in compression can buckle laterally between bracing
    points, often before the truss as a whole approaches its
    design load.
  \item Slender columns in temporary works (formwork supports,
    falsework, scaffolding); these are often designed at the edge
    of code allowable slenderness and have failed historically when
    the slenderness was underestimated or when bracing assumed in
    the design was not provided in execution.
  \item The Quebec Bridge of 1907, in which the lower compression
    chords of the cantilever truss buckled during construction
    under a load below the intended design load; the cause was an
    underestimation of the chord's actual self-weight and an
    overestimate of its buckling capacity from the design tables.
    The collapse killed $75$ workers and is one of the foundational
    structural-engineering case studies. The Royal Commission
    inquiry of 1908 identified errors in the engineer of record's
    buckling calculations as a contributing cause.
  \item Lateral-torsional buckling of unbraced beams under gravity
    load; deep narrow flexural members in tension at the top fibre
    can buckle out of plane in regions where the top fibre is in
    compression (under hogging moments at supports, for example),
    with the failure mode appearing as a sudden lateral deflection
    and twist of the beam.
\end{itemize}

The diagnostic is straightforward: a compression member whose
slenderness exceeds the code limit, whose bracing in service does
not match the design assumption, or whose actual end conditions
differ from the assumed ones, is at risk. The discipline is to
verify the assumptions at fabrication, during installation, and
across the service life. Bracing that exists at design time but is
removed during renovation (an "I'll just take this brace out for a
moment") has caused real collapses.

\subsection{Beam fatigue}

A beam under cyclic loading at amplitudes below the static yield
strength accumulates fatigue damage at locations of stress
concentration. The mechanism follows chapter 3 of this volume; the
specifically beam-related observations are:

\begin{itemize}[leftmargin=*]
  \item Welded steel girder connections accumulate fatigue at the
    weld toe; the toe is the canonical stress concentration in a
    weld and fails preferentially under repeated transverse loads.
  \item Bolted girder splices accumulate fatigue at the net section
    around the bolt holes; the design discipline is to use friction-
    grip bolting that develops sufficient clamping force to avoid
    slip-and-bear cycles.
  \item Highway bridge girders accumulate fatigue from heavy-vehicle
    crossings; the equivalent number of design cycles in a
    fifty-year design life is in the $10^7$ to $10^8$ range for
    typical traffic spectra. Fatigue category curves (Eurocode 3
    chapter 9, AISC 360 Appendix 3) give the stress-cycle limits
    for the various detail types.
  \item Crane girders accumulate fatigue from the trolley's repeated
    crossings; the design discipline is to use fatigue-friendly
    details (rolled section instead of welded, flush welds rather
    than fillet welds at high-stress locations) and to specify
    inspection intervals based on the expected service.
\end{itemize}

% Figure: S-N (Wohler) fatigue curves for welded steel detail categories,
% log-log stress range vs cycles, with the constant-amplitude fatigue
% limit (CAFL) plateau. Illustrative of Eurocode/AISC detail categories.
\begin{figure}[htbp]
\centering
\begin{tikzpicture}
\begin{axis}[
  width=12cm, height=7cm,
  xlabel={number of cycles $N$}, ylabel={stress range $\Delta\sigma$ (MPa)},
  xmode=log, ymode=log,
  xmin=1e4, xmax=1e8, ymin=20, ymax=400,
  axis lines=left,
  legend pos=north east, legend cell align=left,
  legend style={font=\scriptsize},
  tick label style={font=\scriptsize}, label style={font=\small},
  grid=major, grid style={enggray!18},
]
% high-quality detail (category 125): Delta-sigma = 125 at 2e6, slope -1/3
\addplot[engblue, very thick, domain=1e4:5e6, samples=40]
  {125*(2e6/x)^(1/3)};
\addplot[engblue, very thick, domain=5e6:1e8, samples=2]
  coordinates {(5e6,92) (1e8,92)};
\addlegendentry{good detail (cat.\ 125)}
% poor detail (category 50)
\addplot[engred, very thick, domain=1e4:5e6, samples=40]
  {50*(2e6/x)^(1/3)};
\addplot[engred, very thick, domain=5e6:1e8, samples=2]
  coordinates {(5e6,37) (1e8,37)};
\addlegendentry{poor detail (cat.\ 50)}
% reference at 2 million cycles
\addplot[enggray, dashed, thick] coordinates {(2e6,20) (2e6,400)};
\node[enggray, font=\scriptsize, anchor=south, rotate=90] at (axis cs:2e6,55)
  {$N=2\times10^6$ reference};
\node[engblack, font=\scriptsize, anchor=west] at (axis cs:1.2e7,100)
  {fatigue limit (CAFL)};
\end{axis}
\end{tikzpicture}
\caption[S-N curves for welded steel details]{Stress-range versus
cycles-to-failure (S-N, or Wohler) curves for two welded steel detail
categories on log-log axes, in the form used by Eurocode 3 and AISC
360. Each detail category is named by its fatigue strength
$\Delta\sigma_C$ at the reference life of $2\times10^6$ cycles; the
sloping line has slope $-1/3$ in this regime. Below the
constant-amplitude fatigue limit (CAFL plateau), cycles of that
amplitude cause no measurable damage. A smooth, well-ground detail
(blue) tolerates several times the stress range of a sharp,
crack-prone detail (red) at the same life. Detail choice, not base
material, governs the fatigue life of welded structures.}
\label{fig:vol03ch08:fatigue-sn}
\end{figure}


Figure \ref{fig:vol03ch08:fatigue-sn} shows the design tool: the S-N
curve relating the cyclic stress range to the number of cycles a detail
survives, classified by detail category. The detail, not the base
material, governs: a smooth ground transition tolerates several times
the stress range of a sharp re-entrant weld toe at the same life, and
below the constant-amplitude fatigue limit no further damage
accumulates. Fatigue is a slow process; the warning signs (a hairline
crack at a
weld toe, a few millimetres long, visible only on close inspection)
are usually present for months or years before failure. The
discipline of routine inspection is what catches the precursor
before the propagation accelerates to catastrophic.

\subsection{Joint detail failures: cross-reference to Hyatt Regency}

The Hyatt Regency walkway collapse of 17 July 1981, introduced in
chapter 1 of this volume and detailed in the failure section there
\cite{acc:nbs-hyatt-1982}, is the canonical case of a joint detail
failure in this chapter's frame of reference. The design called for
a single continuous threaded rod through the upper and lower
walkways' box-beams; the fabrication-time modification to two
separate rods doubled the load on the upper box-beam's connection,
which then failed by block shear under the doubled load. The
structural engineer of record stamped the modified shop drawings
without recalculating the connection's capacity.

The chapter-8 frame on this case is that the global frame analysis
of the suspended-walkway system would have correctly predicted the
member forces in both the original and the modified designs; what
the analysis did not catch was the local detail failure at the
connection, because the global analysis treats joints as idealised
pins. The discipline that catches detail-level failures is the
\engterm{joint check}: every joint in the structure must be
analysed as its own subassembly, with its own free-body diagram,
its own load paths, and its own failure-mode check. Code-based
design embeds the joint check in the connection-design provisions
(AISC 360 chapter J, Eurocode 3 chapter 8); the engineer who
performs the global frame analysis but skips the joint check has
done half the work.

\subsection{The diagnostic principle for structural failures}

The structural-failure diagnostic question is, at every member and
every joint, which of the failure modes governs:

\begin{itemize}[leftmargin=*]
  \item Material yielding (tension or shear): if the stress at any
    point exceeds the material's yield strength, the member yields
    and may not satisfy serviceability requirements even if it
    does not collapse.
  \item Material fracture (tension): if the stress exceeds the
    ultimate strength, the member fractures.
  \item Buckling (compression): the slenderness ratio and end
    conditions set the buckling capacity; the actual axial load
    must be below it.
  \item Lateral-torsional buckling (bending): the unbraced length
    and section geometry set the critical moment; the actual moment
    must be below it.
  \item Fatigue (cyclic loading): the cyclic-stress amplitude and
    the detail type set the allowable number of cycles; the actual
    service must accumulate fewer than the allowable.
  \item Joint failure (at every connection): the joint's capacity
    against bolt shear, bearing, net-section tension, block shear,
    and fillet-weld fracture must each exceed the joint's actual
    load.
  \item Deflection (serviceability): the actual deflection under
    service loads must be below the code limit.
\end{itemize}

A design that satisfies all seven at every member and joint is a
safe design. A design that fails any of the seven is a vulnerable
design, and the design discipline is to identify which of the seven
will govern and to verify the capacity against the actual demand.

\begin{principle}[Seven structural failure modes, all checked]
A structural design check requires verification against seven
distinct failure modes at every member and joint: yielding,
fracture, buckling, lateral-torsional buckling, fatigue, joint
failure, and deflection. The dominant mode varies by member type
and loading: short columns yield, long columns buckle, deep beams
fatigue at welds, slender beams laterally torsionally buckle, and
every joint can fail at any of several detail-level mechanisms.
The check is mechanical when the modes are enumerated and
catastrophic when one is missed.
\end{principle}

\section{Project}\pathtag{core}

\begin{project}[Design, analyse, build, and verify a small bookshelf]
The reader designs a small wooden bookshelf, performs a full
stress-and-deflection analysis, builds it, and verifies the design
under measured load. The deliverable is a thirty- to fifty-page
engineering report.

\textbf{Track}. Design plus physical build plus analysis.
\textbf{Hazard class}: low.

\textbf{Design}. A free-standing bookshelf, $800$ to $1{,}200\,\si{
\milli\meter}$ wide, $400$ to $800\,\si{\milli\meter}$ tall, with
two to three shelves designed to carry a row of hardcover books
(approximate load $10\,\si{\kilogram}/\text{shelf}/\text{linear}\,
\si{\meter}$). The reader chooses the shelf material (pine,
plywood, MDF), the shelf depth, the support type (end supports,
intermediate bracket, back panel), and the fastener detail.

\textbf{Analysis}. The reader produces:

\begin{itemize}[leftmargin=*]
  \item The free-body diagram of one shelf (chapter 1), the
    reactions at its supports (chapter 2), and the maximum bending
    moment and shear.
  \item The maximum bending stress at the critical section ($\sigma
    = M/S$), compared to the material's allowable bending stress.
  \item The maximum deflection at midspan under service load,
    compared to a code-style limit ($L/360$ or $L/240$ as
    appropriate).
  \item A check of the support columns or end uprights against
    buckling (for shelves taller than they are deep) and against
    yielding of the connection between shelves and uprights.
  \item A check of the bottom-shelf overturning moment if the
    bookshelf is free-standing without wall anchoring.
\end{itemize}

\textbf{Construction}. The reader builds the bookshelf to the
dimensioned design, photographing the construction. Departures from
the design (material substitutions, fastener changes, dimensional
errors) are recorded.

\textbf{Verification}. The reader loads the shelves with their
intended books, measures the deflection at midspan (a ruler is
sufficient), and compares to the predicted deflection.

\textbf{Reflection ($1{,}500$ to $3{,}000$ words)}. The reader
addresses:

\begin{itemize}[leftmargin=*]
  \item Predicted versus measured deflection; reasons for any
    discrepancy.
  \item Which design criterion (stress, deflection, buckling,
    overturning) governed the design.
  \item One feature of the bookshelf whose calculation the reader
    found uncertain; how the uncertainty was handled.
  \item One change to the design the reader would make in a second
    iteration, and the analysis that supports it.
\end{itemize}

\textbf{Deliverable}. The bookshelf (photographed at multiple
stages), the design drawings, the analysis, the verification
measurements, and the reflection.

\textbf{Time}. Four to eight hours of design and analysis, four to
twelve hours of construction, one to two hours of verification, four
to eight hours of write-up.
\end{project}

\section{Exercises}

The exercises mix calculation, derivation, estimation, simulation,
design, diagnosis, failure analysis, and judgment.

\subsection*{Calculation}\pathtag{core}

\begin{exercise}
A simply-supported steel beam ($E = 200\,\text{GPa}$) of span $5\,
\si{\meter}$ has rectangular section $100\,\si{\milli\meter} \times
200\,\si{\milli\meter}$ (depth $200\,\si{\milli\meter}$). It carries
a uniform line load of $5\,\text{kN}/\text{m}$. Find the maximum
bending stress, the maximum deflection, and the section's
slenderness ratio.

\paragraph{Solution sketch.} $M_{\max} = w_0L^2/8 = 5000\times 5^2/8 =
15.6\,\text{kN}\cdot\text{m}$. For the rectangle, $S = bh^2/6 =
0.1\times 0.2^2/6 = 6.67\times 10^{-4}\,\si{\meter\cubed}$ and $I =
bh^3/12 = 6.67\times 10^{-5}\,\si{\meter\tothe{4}}$. Bending stress
$\sigma = M/S = 15\,600/6.67\times 10^{-4} = 23.4\,\text{MPa}$.
Deflection $\delta = 5w_0L^4/(384EI) = 5\times 5000\times 5^4/(384
\times 200\times 10^9 \times 6.67\times 10^{-5}) = 3.05\,\si{\milli
\meter}$ ($\approx L/1640$, well inside any limit). Radius of gyration
$r = \sqrt{I/A} = \sqrt{6.67\times 10^{-5}/0.02} = 57.7\,\si{\milli
\meter}$; treated as a simply supported member of length $5\,\si{\meter}$
the slenderness $L/r = 5000/57.7 \approx 87$.
\end{exercise}

\begin{exercise}
A cantilever beam of length $2\,\si{\meter}$ has a tip point load
of $3\,\text{kN}$. The beam is a steel I-section with $S = 6 \times
10^{-5}\,\si{\meter\cubed}$ (section modulus) and $I = 1.5 \times
10^{-5}\,\si{\meter\tothe{4}}$, $E = 200\,\text{GPa}$. Find the
fixed-end bending moment, the maximum bending stress, and the tip
deflection.

\paragraph{Solution sketch.} Fixed-end moment $M = PL = 3000\times 2 =
6\,\text{kN}\cdot\text{m}$. Bending stress $\sigma = M/S = 6000/(6\times
10^{-5}) = 100\,\text{MPa}$. Tip deflection $\delta = PL^3/(3EI) =
3000\times 2^3/(3\times 200\times 10^9 \times 1.5\times 10^{-5}) =
2.67\,\si{\milli\meter}$.
\end{exercise}

\begin{exercise}
A wooden joist $50 \times 250\,\si{\milli\meter}$ (depth $250\,\si{
\milli\meter}$) supports a uniform residential live load of $2.5\,
\text{kN}/\text{m}$ on a simply-supported $4\,\si{\meter}$ span.
With pine $E = 10\,\text{GPa}$, find the maximum bending stress and
the deflection. Compare to a typical allowable bending stress
($10\,\text{MPa}$) and a deflection limit of $L/360$.

\paragraph{Solution sketch.} $M_{\max} = w_0L^2/8 = 2500\times 4^2/8 =
5\,\text{kN}\cdot\text{m}$. $S = bh^2/6 = 0.05\times 0.25^2/6 =
5.21\times 10^{-4}\,\si{\meter\cubed}$, $I = bh^3/12 = 6.51\times
10^{-5}\,\si{\meter\tothe{4}}$. Bending stress $\sigma = M/S = 5000/
5.21\times 10^{-4} = 9.6\,\text{MPa}$, just under the $10\,\text{MPa}$
allowable. Deflection $\delta = 5w_0L^4/(384EI) = 5\times 2500\times
4^4/(384\times 10\times 10^9 \times 6.51\times 10^{-5}) = 16.0\,\si{
\milli\meter}$, against the limit $L/360 = 4000/360 = 11.1\,\si{\milli
\meter}$. Deflection governs and the joist fails the serviceability
check; a deeper section is required.
\end{exercise}

\begin{exercise}
A steel column of cross-section $W12 \times 65$ (radius of gyration
about the weak axis $r_y \approx 76\,\si{\milli\meter}$, $E = 200\,
\text{GPa}$, $A = 1.23 \times 10^{-2}\,\si{\meter\squared}$) is
pinned at both ends with unbraced length $6\,\si{\meter}$. Find the
Euler critical load.

\paragraph{Solution sketch.} The weak axis governs. $I_y = A r_y^2 =
1.23\times 10^{-2} \times 0.076^2 = 7.10\times 10^{-5}\,\si{\meter
\tothe{4}}$. With $K=1.0$, $P_{\text{cr}} = \pi^2 E I_y/(KL)^2 =
\pi^2\times 200\times 10^9 \times 7.10\times 10^{-5}/6^2 = 3.89\times
10^6\,\text{N} \approx 3890\,\text{kN}$. Check slenderness $\lambda =
KL/r_y = 6000/76 = 79 < \lambda_c \approx 89$, so the column is
intermediate and the elastic Euler load overstates the real capacity.
\end{exercise}

\begin{exercise}
A steel pipe column of outer diameter $100\,\si{\milli\meter}$,
wall thickness $4\,\si{\milli\meter}$, length $3\,\si{\meter}$, fixed
at one end and free at the other, carries an axial compressive load.
Find the Euler critical load.

\paragraph{Solution sketch.} Thin tube, $d_o = 0.100$, $d_i = 0.092$.
$I = \pi(d_o^4 - d_i^4)/64 = \pi(0.100^4 - 0.092^4)/64 = 1.38\times
10^{-6}\,\si{\meter\tothe{4}}$. Fixed-free gives $K = 2.0$, so
$P_{\text{cr}} = \pi^2 E I/(KL)^2 = \pi^2\times 200\times 10^9 \times
1.38\times 10^{-6}/(2\times 3)^2 = 7.57\times 10^4\,\text{N} \approx
75.7\,\text{kN}$. The $K=2$ flagpole condition cuts the capacity to a
quarter of the pinned-pinned value.
\end{exercise}

\begin{exercise}
A steel cantilever of length $1\,\si{\meter}$ has a tip load
$P$ that produces a bending stress of $200\,\text{MPa}$ at the
fixed end. The cantilever has a rectangular cross-section $30
\times 50\,\si{\milli\meter}$ (depth $50\,\si{\milli\meter}$). Find
$P$.

\paragraph{Solution sketch.} $S = bh^2/6 = 0.030\times 0.050^2/6 =
1.25\times 10^{-5}\,\si{\meter\cubed}$. Fixed-end moment $M = \sigma S
= 200\times 10^6 \times 1.25\times 10^{-5} = 2500\,\text{N}\cdot\text{m}$.
For a cantilever $M = PL$, so $P = M/L = 2500/1 = 2500\,\text{N} =
2.5\,\text{kN}$.
\end{exercise}

\begin{exercise}
A circular shaft of $40\,\si{\milli\meter}$ diameter is subjected to
combined torsion $T = 500\,\text{N}\cdot\text{m}$ and bending moment
$M = 600\,\text{N}\cdot\text{m}$. Find the maximum principal stress
and the maximum shear stress at the shaft's surface.

\paragraph{Solution sketch.} This is example 3 of the worked-examples
section with $d = 0.040\,\si{\meter}$. $\pi d^3/32 = 6.28\times 10^{-6}
\,\si{\meter\cubed}$, giving $\sigma = M/(\pi d^3/32) = 95.5\,\text{MPa}$
and $\tau = T/(\pi d^3/16) = 39.8\,\text{MPa}$. The principal stresses
$\sigma_{1,2} = \sigma/2 \pm \sqrt{(\sigma/2)^2 + \tau^2} = 47.7 \pm
62.1\,\text{MPa}$ give $\sigma_1 = 109.8\,\text{MPa}$ and $\tau_{\max}
= 62.1\,\text{MPa}$.
\end{exercise}

\begin{exercise}
An indeterminate beam (both ends fixed) of span $L = 6\,\si{\meter}$
carries a centred point load $P = 20\,\text{kN}$. Find the
fixed-end moments and the midspan moment.

\paragraph{Solution sketch.} For a fixed-fixed beam with a central
point load, the end moments are $M_{\text{end}} = PL/8$ (hogging) and
the midspan moment is $M_{\text{mid}} = PL/8$ (sagging). With $P =
20\,\text{kN}$, $L = 6\,\si{\meter}$: $M_{\text{end}} = M_{\text{mid}} =
20\times 6/8 = 15\,\text{kN}\cdot\text{m}$. The fixed ends halve the
peak moment relative to the simply supported $PL/4 = 30\,\text{kN}\cdot
\text{m}$, the value of end fixity.
\end{exercise}

\begin{exercise}
Use Castigliano's theorem to find the tip deflection of a cantilever
beam of length $L$ under a tip load $P$. Verify that the result
agrees with $\delta = PL^3/(3EI)$.

\paragraph{Solution sketch.} Measure $x$ from the tip; the moment at
$x$ is $M(x) = -Px$. The strain energy is $U = \int_0^L M^2/(2EI)\,dx
= \int_0^L P^2x^2/(2EI)\,dx = P^2L^3/(6EI)$. By Castigliano,
$\delta = \partial U/\partial P = 2P L^3/(6EI) = PL^3/(3EI)$, which
matches the table value. Equivalently the unit-load method with
$m(x) = -x$ gives $\delta = \int_0^L M(x)m(x)/(EI)\,dx = \int_0^L
Px^2/(EI)\,dx = PL^3/(3EI)$.
\end{exercise}

\begin{exercise}
A steel column has slenderness ratio $\lambda = KL/r = 120$. The
material has $E = 200\,\text{GPa}$ and $\sigma_y = 250\,\text{MPa}$.
Determine whether the column is short, intermediate, or long, and
compute its Euler buckling stress.

\paragraph{Solution sketch.} Transition slenderness $\lambda_c = \pi
\sqrt{E/\sigma_y} = \pi\sqrt{200\times 10^9/250\times 10^6} = \pi
\sqrt{800} = 88.9$. Since $\lambda = 120 > \lambda_c$, the column is
long and governed by elastic Euler buckling. The Euler stress is
$\sigma_{\text{cr}} = \pi^2E/\lambda^2 = \pi^2\times 200\times 10^9/
120^2 = 137\,\text{MPa}$, well below $\sigma_y = 250\,\text{MPa}$,
confirming buckling rather than yielding governs.
\end{exercise}

\subsection*{Derivation}\pathtag{standard}

\begin{exercise}
Derive the Euler-Bernoulli beam equation $EI \,y'''' = w(x)$ from
the moment-curvature relation, the moment-shear-load relations, and
the geometric small-deflection identity $\kappa \approx y''$.

\paragraph{Solution sketch.} Start from the moment-curvature relation
$M = EI\kappa = EI\,y''$. Differentiate twice with respect to $x$,
holding $EI$ constant: $M'' = EI\,y''''$. The equilibrium relations are
$dM/dx = V$ and $dV/dx = -w$, so $M'' = (dM/dx)' = V' = dV/dx = -w$ (or
$+w$ depending on the chosen sign convention for $w$). Substituting,
$EI\,y'''' = w(x)$. The chain is moment-curvature ($M$ to $y''$), then
two derivatives carry $M$ down to the load through the two equilibrium
relations.
\end{exercise}

\begin{exercise}
Derive the maximum-deflection formula $\delta = 5 w_0 L^4/(384 EI)$
for a simply-supported beam under uniform line load. State the
boundary conditions used.

\paragraph{Solution sketch.} Integrate $EI\,y'''' = w_0$ four times:
$EI\,y''' = w_0 x + C_1$, $EI\,y'' = w_0 x^2/2 + C_1 x + C_2$, and so
on. Boundary conditions for simple supports at $x = 0, L$: $y = 0$ and
$M = EI\,y'' = 0$ at each end. The two moment conditions give $C_2 = 0$
and $C_1 = -w_0L/2$ (so $y'' = (w_0/2EI)(x^2 - Lx)$). Integrating twice
more with $y(0) = y(L) = 0$ gives $y(x) = (w_0/24EI)(x^4 - 2Lx^3 +
L^3x)$. At midspan $x = L/2$, $y(L/2) = (w_0/24EI)(L^4/16 - L^4/4 +
L^4/2) = (w_0/24EI)(5L^4/16) = 5w_0L^4/(384EI)$.
\end{exercise}

\begin{exercise}
Derive the Euler critical load $P_{\text{cr}} = \pi^2 EI/L^2$ for a
pinned-pinned column from the moment-equilibrium equation $EI y'' +
P y = 0$ and the boundary conditions $y(0) = y(L) = 0$.

\paragraph{Solution sketch.} Write $k^2 = P/EI$, so $y'' + k^2 y = 0$
with general solution $y = A\sin kx + B\cos kx$. The condition $y(0) =
0$ forces $B = 0$. The condition $y(L) = 0$ then requires $A\sin kL =
0$; the non-trivial solution ($A \neq 0$) needs $\sin kL = 0$, hence
$kL = n\pi$ for integer $n$. The smallest non-zero load corresponds to
$n = 1$, $kL = \pi$, so $k^2 = \pi^2/L^2 = P/EI$, giving $P_{\text{cr}}
= \pi^2EI/L^2$. The mode shape is the half-sine $y = A\sin(\pi x/L)$.
\end{exercise}

\begin{exercise}
Show that the effective-length factor for a fixed-fixed column is
$K = 0.5$. Identify the buckled mode shape and its inflection
points.

\paragraph{Solution sketch.} For a fixed-fixed column the buckled shape
is the full cosine wave $y \propto 1 - \cos(2\pi x/L)$, with zero slope
at both ends as the fixed condition requires. The critical load works
out to $P_{\text{cr}} = 4\pi^2EI/L^2 = \pi^2EI/(L/2)^2$, so comparison
with $\pi^2EI/(KL)^2$ gives $K = 0.5$. The inflection points (where
$y'' = 0$, the moment is zero) are at $x = L/4$ and $x = 3L/4$; the
distance between them is $L/2$, which is the effective length $KL$.
Between these inflection points the column behaves like a pinned-pinned
column of length $L/2$.
\end{exercise}

\begin{exercise}
Derive Castigliano's first theorem $\delta_i = \partial U/\partial
P_i$ for a linear elastic structure. State the assumption that makes
the strain energy depend only on the loads and not on the loading
history.

\paragraph{Solution sketch.} For a linear elastic structure the
displacements are linear in the loads, so the work done is $U =
\tfrac12 \sum_j P_j \delta_j$ and depends only on the final load set,
not the order of application (the load-path independence assumption,
valid because the strain energy is a state function for an elastic
body). Increase load $P_i$ by $dP_i$; the additional work is
$\delta_i\,dP_i$ to first order (the new load times the displacement
already present at its own point), so $\partial U/\partial P_i =
\delta_i$. The reciprocal theorem (Maxwell-Betti), which follows from
the symmetry of the flexibility matrix, underwrites the equality.
\end{exercise}

\begin{exercise}
Show that, for a slender beam under bending and axial load, the
P-delta amplification factor for the first-order deflection is
approximately $1/(1 - P/P_{\text{cr}})$ for $P < P_{\text{cr}}$.

\paragraph{Solution sketch.} Let $\delta_0$ be the first-order
(transverse-load only) deflection and $\delta$ the amplified
deflection. The axial force $P$ acting through the deflection adds a
moment $P\delta$, and the extra deflection it causes is proportional to
that moment. Expanding the deflection as a sine series and keeping the
first mode, each P-delta cycle multiplies the contribution by
$P/P_{\text{cr}}$, giving the geometric series $\delta = \delta_0(1 +
P/P_{\text{cr}} + (P/P_{\text{cr}})^2 + \cdots) = \delta_0/(1 -
P/P_{\text{cr}})$. The amplification is mild for $P \ll P_{\text{cr}}$
and diverges as $P \to P_{\text{cr}}$, which is the same instability the
Euler load marks.
\end{exercise}

\subsection*{Estimation}\pathtag{core}

\begin{exercise}
Estimate the maximum span of a pine $50 \times 150\,\si{\milli\meter}$
floor joist (depth $150\,\si{\milli\meter}$) carrying a residential
live load of $2\,\text{kN}/\text{m}$. State which design criterion
(bending stress, deflection) governs.

\paragraph{Reference estimate.} Pine: $E \approx 10\,\text{GPa}$,
$\sigma_{\text{allow}} \approx 10\,\text{MPa}$. Section $50\times
150\,\si{\milli\meter}$: $S = bh^2/6 = 0.05\times 0.15^2/6 = 1.875
\times 10^{-4}\,\si{\meter\cubed}$, $I = bh^3/12 = 1.41\times 10^{-5}
\,\si{\meter\tothe{4}}$. Bending limit: $L^2 = 8S\sigma/w_0 = 8\times
1.875\times 10^{-4}\times 10^7/2000 = 7.5$, $L \approx 2.74\,\si{\meter}$.
Deflection limit $L/360$: $L^3 = 384EI/(1800 w_0) = 384\times 10^{10}
\times 1.41\times 10^{-5}/(1800\times 2000) = 15.0$, $L \approx 2.47\,
\si{\meter}$. Deflection governs at $L \approx 2.5\,\si{\meter}$, the
usual outcome for shallow residential joists.
\end{exercise}

\begin{exercise}
Estimate the buckling load of a wooden $25 \times 25\,\si{\milli\meter}$
column of $1\,\si{\meter}$ length, with one end fixed and one end
free. State the assumed material properties and the section's
radius of gyration.

\paragraph{Reference estimate.} Wood $E \approx 10\,\text{GPa}$.
Section $25\times 25\,\si{\milli\meter}$: $I = bh^3/12 = 0.025^4/12 =
3.26\times 10^{-8}\,\si{\meter\tothe{4}}$, $A = 6.25\times 10^{-4}\,
\si{\meter\squared}$, $r = \sqrt{I/A} = 7.2\,\si{\milli\meter}$.
Fixed-free $K = 2$, $L = 1\,\si{\meter}$, so $P_{\text{cr}} = \pi^2EI/
(KL)^2 = \pi^2\times 10\times 10^9 \times 3.26\times 10^{-8}/(2)^2 =
805\,\text{N}$, roughly $80\,\si{\kilogram}$ of axial load. Slenderness
$KL/r = 2000/7.2 \approx 278$, very slender, so elastic buckling
firmly governs.
\end{exercise}

\begin{exercise}
Estimate the tip deflection of a typical broom handle of $1\,\si{
\meter}$ length and $20\,\si{\milli\meter}$ outer diameter under a
$5\,\si{\newton}$ side push at its tip. State assumptions about
material and cross-section.

\paragraph{Reference estimate.} Treat as a cantilever, wood $E \approx
10\,\text{GPa}$. Solid circle $d = 20\,\si{\milli\meter}$: $I = \pi
d^4/64 = \pi\times 0.020^4/64 = 7.85\times 10^{-9}\,\si{\meter\tothe{4}}$.
Tip deflection $\delta = PL^3/(3EI) = 5\times 1^3/(3\times 10\times
10^9 \times 7.85\times 10^{-9}) = 0.021\,\si{\meter} \approx 21\,\si{
\milli\meter}$. A broom handle visibly bends a couple of centimetres
under a modest side push, which matches everyday experience; a hollow
handle would deflect somewhat more.
\end{exercise}

\begin{exercise}
Estimate the bending stress at the base of a $5\,\si{\meter}$
street lamp pole in a $20\,\si{\meter\per\second}$ wind. State
assumptions about pole geometry, wind drag, and material.

\paragraph{Reference estimate.} Assume a steel tube, outer diameter
$150\,\si{\milli\meter}$, wall $5\,\si{\milli\meter}$, height $5\,\si{
\meter}$, projected width $0.15\,\si{\meter}$. Dynamic pressure $q =
\tfrac12\rho v^2 = 0.5\times 1.2\times 20^2 = 240\,\text{Pa}$. With a
drag coefficient $C_d \approx 1.2$ for a cylinder, the line load is
$w = C_d q D = 1.2\times 240\times 0.15 = 43\,\text{N}/\text{m}$. As a
cantilever the base moment is $M = wH^2/2 = 43\times 5^2/2 = 540\,
\text{N}\cdot\text{m}$. Tube $I = \pi(0.150^4 - 0.140^4)/64 = 6.2
\times 10^{-6}\,\si{\meter\tothe{4}}$, $S = I/(D/2) = 8.3\times 10^{-5}
\,\si{\meter\cubed}$. Base bending stress $\sigma = M/S = 540/8.3
\times 10^{-5} \approx 6.5\,\text{MPa}$, modest, so a lamp pole is
controlled by fatigue and gust dynamics rather than static strength.
\end{exercise}

\begin{exercise}
Estimate the maximum unsupported height of a $200\,\si{\milli
\meter}$-square reinforced-concrete column of typical residential
design. State the slenderness limit and the assumed material.

\paragraph{Reference estimate.} A reinforced-concrete column is
treated as ``short'' (no slenderness reduction) below a slenderness
ratio of about $kL_u/r \leq 22$ for a braced frame. For a square
section $b = 200\,\si{\milli\meter}$, $r = b/\sqrt{12} = 200/3.46 =
57.7\,\si{\milli\meter}$. With $k \approx 1.0$ (pinned-pinned) the
limit gives $L_u \leq 22 r = 22\times 0.0577 = 1.27\,\si{\meter}$;
with end fixity ($k \approx 0.7$) the height roughly extends to
$\sim 1.8\,\si{\meter}$. A storey-height column ($\sim 3\,\si{\meter}$)
of this small section is therefore slender and needs a moment
magnification check; a residential column is usually larger
($300\,\si{\milli\meter}$ or more) to stay short.
\end{exercise}

\subsection*{Simulation}\pathtag{mastery}

\begin{exercise}
Implement a numerical solver for the Euler-Bernoulli beam equation
$EI\,y''''(x) = w(x)$ using the finite-difference method. Verify on
a simply-supported beam under uniform load against the analytical
solution.

\paragraph{Solution sketch.} The reference implementation is
\path{code/beam_deflection.py}. Discretise the span into $n$ nodes,
approximate $y''''$ by the five-point central difference $(y_{i-2} -
4y_{i-1} + 6y_i - 4y_{i+1} + y_{i+2})/h^4$, assemble the banded linear
system, and impose $y = 0$ and $y'' = 0$ (a ghost-node reflection) at
the simple supports. Solving gives the midspan deflection within about
$0.1\,\%$ of $5w_0L^4/(384EI)$ on a two-hundred-node grid; the error
falls as $h^2$ as the grid is refined, the signature of the
second-order difference. The script writes the profile to
\path{data/deflection_profile.csv}.
\end{exercise}

\begin{exercise}
Implement the stiffness method for a planar frame: input member
connectivity, member properties ($EA$, $EI$, length), joint
coordinates, applied loads, support types; output joint
displacements and member end-forces. Verify on a textbook portal
frame.

\paragraph{Solution sketch.} The reference implementation is
\path{code/frame_stiffness.py}. Build the $6\times 6$ local member
matrix (axial $EA/L$ terms coupled with the flexural $12EI/L^3$,
$6EI/L^2$, $4EI/L$, $2EI/L$ terms), rotate it to global axes with the
direction-cosine transformation $T^{\!\top}kT$, and assemble member
contributions into the global $3n\times 3n$ matrix by the connectivity.
Delete the rows and columns of the restrained degrees of freedom and
solve $K\vec{u} = \vec{F}$. The script checks the joint rotation of a
propped cantilever against the closed-form $ML/(4EI)$ and agrees to
better than $0.1\,\%$; verification on a portal frame compares joint
sway and member end-moments to a published worked example or to a
moment-distribution hand check.
\end{exercise}

\begin{exercise}
Simulate the post-buckling behaviour of a pinned-pinned column
beyond its Euler critical load, including the geometric
nonlinearity (large deflections). Plot the load versus midspan
deflection and identify the supercritical pitchfork bifurcation.

\paragraph{Solution sketch.} Use the elastica (large-deflection)
formulation, in which the curvature is $d\theta/ds$ along the arc
length $s$ rather than $y''$, giving $EI\,d^2\theta/ds^2 + P\sin\theta
= 0$, the pendulum equation in disguise. Solve the two-point boundary
value problem (shooting on the end slope) for a sequence of loads
$P > P_{\text{cr}}$ and record the midspan lateral deflection. The
load-deflection curve is flat (zero deflection) up to $P_{\text{cr}} =
\pi^2EI/L^2$, then rises along the two symmetric branches sketched in
figure \ref{fig:vol03ch08:postbuckling}: the supercritical pitchfork.
The deflected branch follows $P/P_{\text{cr}} \approx 1 + (\pi^2/8)
(\delta/L)^2$ near the bifurcation. Adding a small initial
imperfection rounds the corner and removes the sharp branch point, as
real columns show.
\end{exercise}

\subsection*{Design}\pathtag{standard}

\begin{exercise}
Design a simply-supported steel beam to carry a uniform live load
of $10\,\text{kN}/\text{m}$ over a $6\,\si{\meter}$ span, using
$A992$ structural steel with $\sigma_y = 345\,\text{MPa}$. Specify
the section (from a standard rolled-section table) and check
against bending stress, shear stress, and a deflection limit of
$L/360$.

\paragraph{Model-answer sketch.} $M_{\max} = w_0L^2/8 = 10\,000\times
6^2/8 = 45\,\text{kN}\cdot\text{m}$. With $\phi M_n = \phi\sigma_y S$
and $\phi = 0.9$, the required $S = M/(\phi\sigma_y) = 45\,000/(0.9
\times 345\times 10^6) = 1.45\times 10^{-4}\,\si{\meter\cubed} =
145\,\text{cm}^3$. Deflection: $I \geq 360\times 5 w_0 L^3/(384E) =
1800\times 10\,000\times 6^3/(384\times 200\times 10^9) = 5.06\times
10^{-5}\,\si{\meter\tothe{4}} = 5060\,\text{cm}^4$. From
\path{data/standard_sections.csv}, a $W16\times 40$ ($S = 1070\,\text{cm}^3$,
$I = 21\,800\,\text{cm}^4$) passes both with margin; check web shear
$\tau = V/A_{\text{web}}$ with $V = w_0L/2 = 30\,\text{kN}$, easily
satisfied. A lighter $W12\times 65$ also works but is heavier than the
deeper $W16$, which is the economical pick.
\end{exercise}

\begin{exercise}
Design a steel column to support an axial load of $300\,\text{kN}$
over an unbraced length of $4\,\si{\meter}$, with pinned ends. Use
$A992$ steel ($\sigma_y = 345\,\text{MPa}$). Specify the section
and check against both Euler buckling and material yielding.

\paragraph{Model-answer sketch.} Try a $W8\times 31$ ($A = 58.7\,\text{cm}^2$,
weak-axis $r_y = 51\,\si{\milli\meter}$). Slenderness $\lambda = KL/r_y
= 1.0\times 4000/51 = 78 < \lambda_c \approx 76$ to $89$ for $A992$
($\lambda_c = \pi\sqrt{E/\sigma_y} = \pi\sqrt{200\times 10^9/345\times
10^6} = 75.6$), so the column is just into the long range. Euler stress
$\sigma_{\text{cr}} = \pi^2E/\lambda^2 = \pi^2\times 200\times 10^9/78^2
= 324\,\text{MPa}$; the inelastic AISC column equation reduces this to
about $0.7\sigma_y \approx 240\,\text{MPa}$. Capacity $\phi P_n \approx
0.9\times 240\times 10^6 \times 58.7\times 10^{-4} = 1270\,\text{kN}
\gg 300\,\text{kN}$ demand. The squash load $\sigma_y A = 2025\,\text{kN}$
confirms yielding does not govern. The $W8\times 31$ is generously
sufficient; a smaller section could be tried if economy matters.
\end{exercise}

\begin{exercise}
Design a glulam timber beam to carry a residential floor live load
of $2.5\,\text{kN}/\text{m}$ over a $6\,\si{\meter}$ span, with
$E = 12\,\text{GPa}$ and $\sigma_{\text{allow}} = 12\,\text{MPa}$.
Specify the cross-section and check both bending and deflection.

\paragraph{Model-answer sketch.} $M_{\max} = w_0L^2/8 = 2500\times 6^2/
8 = 11.25\,\text{kN}\cdot\text{m}$. Bending: $S \geq M/\sigma_{\text{allow}}
= 11\,250/12\times 10^6 = 9.4\times 10^{-4}\,\si{\meter\cubed} =
940\,\text{cm}^3$. Deflection $L/360$: $I \geq 1800 w_0 L^3/(384E) =
1800\times 2500\times 6^3/(384\times 12\times 10^9) = 2.11\times 10^{-4}
\,\si{\meter\tothe{4}}$. Try a glulam $130\times 400\,\si{\milli\meter}$:
$S = bh^2/6 = 0.13\times 0.4^2/6 = 3.47\times 10^{-3}\,\si{\meter\cubed}$
(ample) and $I = bh^3/12 = 6.93\times 10^{-4}\,\si{\meter\tothe{4}}$
(ample). Deflection governs the depth; a $130\times 360\,\si{\milli
\meter}$ section, $I = 5.05\times 10^{-4}\,\si{\meter\tothe{4}}$, still
clears the $2.11\times 10^{-4}$ requirement, so the $360\,\si{\milli
\meter}$ depth is the economical choice (\path{data/timber_design_values.csv}
for the glulam grade).
\end{exercise}

\begin{exercise}
Design a moment-resisting steel portal frame for a $10\,\si{\meter}$
span and $4\,\si{\meter}$ eaves height, carrying a uniform roof
live load of $3\,\text{kN}/\text{m}$. Specify the column and rafter
sections and the moment connection at the eaves.

\paragraph{Model-answer sketch.} For a pinned-base portal under gravity
load $w_0$ over span $L$, the eaves moment is roughly $M_{\text{eaves}}
\approx w_0L^2/16 = 3000\times 10^2/16 = 18.75\,\text{kN}\cdot\text{m}$
and the ridge sagging moment $\sim w_0L^2/24$; a frame analysis with
\path{code/frame_stiffness.py} sharpens these. Size the rafter for the
larger of the eaves and span moments: required $S \approx M/(\phi\sigma_y)
= 18\,750/(0.9\times 345\times 10^6) = 6.0\times 10^{-5}\,\si{\meter
\cubed} = 60\,\text{cm}^3$, a modest section such as a $W8\times 31$.
Columns carry the eaves moment plus axial; check combined action with
the interaction equation $N/N_{\text{allow}} + M/M_{\text{allow}} \leq
1$. The eaves moment connection must transfer the full $18.75\,\text{kN}
\cdot\text{m}$ through a haunched or fully welded rigid joint; specify a
bolted end-plate or welded haunch detailed per the connection code.
\end{exercise}

\begin{exercise}
Design a cantilever balcony slab of $1.5\,\si{\meter}$ projection
to support a residential live load of $2.5\,\text{kN}/\text{m}^2$
plus its own self-weight. Specify the slab thickness, the
reinforcement (or equivalent if not reinforced concrete), and the
deflection criterion.

\paragraph{Model-answer sketch.} Per metre width, the load is $w =
(2.5 + \text{self weight})\,\text{kN}/\text{m}^2$. For a $180\,\si{
\milli\meter}$ slab the self-weight is $0.18\times 24 = 4.3\,\text{kN}/
\text{m}^2$, so $w \approx 6.8\,\text{kN}/\text{m}$ per metre strip.
Cantilever moment at the support $M = wL^2/2 = 6.8\times 1.5^2/2 =
7.65\,\text{kN}\cdot\text{m}$ per metre. This tension-on-top moment is
resisted by top reinforcement; sizing follows the concrete-design
chapter (Vol V), giving on the order of a $\varnothing12$ bar at
$150\,\si{\milli\meter}$ centres in the top face. The governing
serviceability check is the cantilever deflection limit, typically
$L/180$ for a cantilever ($1500/180 = 8.3\,\si{\milli\meter}$), which
sets the slab thickness; codes give a span-to-depth ratio near $7$ for
cantilevers, so $180\,\si{\milli\meter}$ for a $1.5\,\si{\meter}$
projection is reasonable. Detail the top steel to extend a full
development length back into the supporting slab.
\end{exercise}

\subsection*{Diagnosis and reverse engineering}\pathtag{core}

\begin{exercise}
A residential floor visibly bounces when occupants walk across it.
Identify the most probable cause, the diagnostic measurement, and
the corrective action.

\paragraph{Solution sketch.} The probable cause is inadequate floor
stiffness, not inadequate strength: the joists satisfy the bending
check but the deflection (or the natural frequency, which scales with
stiffness) is low enough that footfall excites a perceptible response.
Diagnostic measurement: deflect the floor with a known weight at
midspan and compare to the predicted $\delta$, or measure the floor's
fundamental frequency (perceptible bounce typically appears below
$\sim 8\,\text{Hz}$). Corrective action: increase stiffness, by adding
a midspan beam to halve the span, sistering deeper joists, adding
bridging or blocking to share load between joists, or stiffening the
deck. Adding stiffness, not strength, is the fix, because $EI$ governs
both deflection and frequency.
\end{exercise}

\begin{exercise}
A steel column in a parking structure has been observed to deflect
visibly under design vehicle loading. Identify the failure mode at
risk, the inspection priority, and the corrective action.

\paragraph{Solution sketch.} A column that deflects laterally under
axial load is showing P-delta behaviour: the visible lateral movement
means the axial load is an appreciable fraction of the buckling load,
and the deflection amplifies the moment in the feedback described in
section 8.3. The failure mode at risk is column buckling (instability),
potentially sudden. Inspection priority: verify the actual end
conditions and bracing against the design assumptions, check for an
overload (heavier vehicles than designed, or an added storey), and look
for any out-of-plumb or local damage that raises the effective
slenderness. Corrective action: reduce the effective length by adding
lateral bracing at mid-height (which can multiply the buckling capacity
several-fold), reduce the load, or strengthen the section. Treat this as
urgent, since buckling gives little warning.
\end{exercise}

\begin{exercise}
A welded steel girder bridge has developed a small fatigue crack at
a stiffener-to-flange weld toe after twenty years of service.
Identify the design discipline that should have anticipated this
failure mode and the inspection regime that should have caught the
crack earlier.

\paragraph{Solution sketch.} The failure mode is high-cycle fatigue at
a weld-toe stress concentration; the stiffener-to-flange weld toe is a
classic poor fatigue detail (a low detail category on the S-N curve of
figure \ref{fig:vol03ch08:fatigue-sn}). The design discipline that
should have anticipated it is fatigue design: classify each welded
detail by category, compute the stress range from the live-load
spectrum, and verify the accumulated damage against the category's S-N
curve over the design life (Eurocode 3 part 1-9, AISC 360 Appendix 3).
The inspection regime: scheduled visual and dye-penetrant or
magnetic-particle inspection of fatigue-critical details at intervals
set by the predicted crack-growth life, so a millimetre-scale toe crack
is found and ground out or retrofitted long before it propagates
through the flange.
\end{exercise}

\begin{exercise}
A wooden bookcase has tipped forward under load. The bookcase was
free-standing without wall anchoring. Identify the failure mode and
the design correction.

\paragraph{Solution sketch.} The failure mode is overturning
(rigid-body instability), not a strength failure of any member. The
bookcase tips when the overturning moment of the load (books on the
upper shelves, or a horizontal pull) about the front edge of the base
exceeds the restoring moment of the self-weight acting through the
centre of gravity. The relevant inequality is $W\cdot(b/2) \geq
F_h\cdot h$, restoring versus overturning. Design corrections: anchor
the top to the wall (the standard fix, and code-required in many
jurisdictions for tall furniture), widen or weight the base to move the
centre of gravity forward and lower it, or keep heavy items on the
lower shelves. Anchoring is the reliable correction because it removes
the dependence on loading pattern.
\end{exercise}

\subsection*{Failure analysis}\pathtag{core}

\begin{exercise}
Re-read the NBS report on the Hyatt Regency walkway collapse \cite{
acc:nbs-hyatt-1982}. Identify the seven structural failure modes
(section 8.8) and state which one was the proximate cause of the
collapse. State the design-time analysis that should have caught
the modification's effect on that failure mode.

\paragraph{Solution sketch.} Of the seven modes (yielding, fracture,
buckling, lateral-torsional buckling, fatigue, joint failure,
deflection), the proximate cause was joint failure: the box-beam
connection to the support rod failed under load. The fabrication change
from one continuous rod to two stacked rods doubled the load on the
upper-walkway connection, which then failed as the nut tore through the
box-beam (a block-shear / weld-tearout joint mechanism). The
design-time analysis that should have caught it is the joint check: a
free-body diagram of the connection treated as its own subassembly,
recomputed for the modified detail when the shop drawings changed. The
global frame forces were unchanged; the local connection capacity was
halved, and only a per-joint check exposes that. See
\cite{acc:nbs-hyatt-1982}.
\end{exercise}

\begin{exercise}
A long-span steel truss bridge has buckled one of its compression
chords during construction. Argue, in $200$ to $300$ words, what
the design-time slenderness check should have produced, what the
construction-time bracing should have provided, and what
construction discipline was missing.

\paragraph{Solution sketch.} The design-time slenderness check should
have established the chord's effective length between brace points and
verified $\lambda = KL/r$ against the code column curve for the actual
construction-stage load, including the chord's self-weight (the Quebec
Bridge case turned on an underestimate of exactly this). A compression
chord is only as strong as its unbraced length allows; if the bracing
assumed in design is not yet installed during erection, the effective
length is longer and the buckling load far lower. The construction-time
bracing should have provided lateral support at the spacing the design
assumed, present before the chord was loaded. The missing discipline is
construction-stage analysis: verifying that the partially erected
structure is stable at every stage, with the bracing actually in place
matching the bracing the strength calculation assumed. A strong
finished structure can collapse mid-erection if a temporary
configuration is unbraced.
\end{exercise}

\begin{exercise}
A multi-storey building's beam-to-column moment connection has
fractured along the column flange under earthquake loading. State,
in $200$ to $300$ words, the most probable failure mode (brittle
weld fracture, ductile yielding, fastener failure) and the design
discipline that catches this mode (typically post-Northridge
detailing requirements in seismic design).

\paragraph{Solution sketch.} The most probable mode is brittle weld
fracture: a fracture along the column flange under cyclic seismic
demand points to a weld and heat-affected-zone failure before the
intended ductile yielding of the beam could develop. This is the
signature of pre-Northridge welded moment connections, where the
complete-joint-penetration weld at the beam flange, with its backing
bars and low-toughness weld metal, fractured at low ductility. The
design discipline that catches it is capacity (ductile) seismic
detailing: force the inelastic action into a pre-qualified ductile
fuse (a reduced beam section, or a connection detailed to keep yielding
in the beam away from the weld), use tough weld metal, remove or treat
backing bars, and qualify the connection by cyclic testing. The
post-Northridge connection provisions exist precisely to move the
failure from brittle weld fracture to ductile beam yielding.
\end{exercise}

\subsection*{Judgment}\pathtag{standard}

\begin{exercise}
A junior engineer asks whether the deflection limit $L/360$ is
based on safety or on serviceability. Write a one-paragraph reply
that takes the question seriously and identifies what the limit
actually protects.

\paragraph{Solution sketch.} The $L/360$ limit is serviceability, not
safety: a beam can deflect well past $L/360$ without any risk of
collapse, since strength is checked separately and with its own safety
factors. What the limit protects is everything attached to the beam
that cannot tolerate movement: plastered ceilings and partitions that
crack, doors and windows that bind, finishes that show distress, and
occupants who perceive a floor as too flexible. The specific ratio is
calibrated to the supported elements, which is why brittle finishes
(glazing, masonry) get a tighter limit such as $L/500$ or $L/600$ while
a roof with no ceiling tolerates $L/180$. A good reply distinguishes the
strength limit state from the serviceability limit state and names what
the deflection limit is actually defending.
\end{exercise}

\begin{exercise}
A senior colleague argues that hand-calculation methods (moment
distribution, slope-deflection, unit-load) are obsolete now that
finite-element software runs on every laptop. Write a one-paragraph
reply that states when hand methods are still appropriate.

\paragraph{Solution sketch.} Hand methods remain the right tool in
several situations even with universal software. They give the
order-of-magnitude check that catches a gross input error in a
finite-element model (a misplaced support, a unit slip, a load applied
to the wrong node); an engineer who cannot estimate the answer by hand
cannot tell whether the software output is sensible. They are faster
than building a model for small or preliminary problems, where a
moment-distribution or unit-load calculation answers the question in
minutes. They build the intuition for load paths that lets an engineer
read a model critically rather than trusting it. A good reply concedes
that software is the production tool for large analyses while insisting
that hand methods are the verification and intuition layer that keeps
the software honest.
\end{exercise}

\begin{exercise}
A designer proposes to remove a column from an existing building to
create an open-plan office. Write a one-paragraph response covering
the structural-analysis steps that must be completed and the
verification stages required before the modification can be
approved.

\paragraph{Solution sketch.} Removing a column changes the load path:
the loads the column carried must be rerouted to the remaining
structure, usually by a new transfer beam spanning the opening. The
analysis steps are to trace what the column supports (the tributary
gravity load and any role in the lateral system), establish the new
load path, design the transfer beam and its connections for the
redistributed forces, and check the remaining columns and foundations
for the increased load. The deflection of the new transfer beam must be
limited so the structure above is not distressed. Verification stages:
an independent design check, a temporary shoring scheme to carry the
load during the work, staged removal with monitoring, and a final
inspection confirming the as-built matches the design. The response
should stress that this is a primary structural alteration requiring a
qualified engineer, a permit, and shoring, not a casual change.
\end{exercise}

\begin{exercise}
Identify a beam or column in a building you regularly use whose
load path you have never explicitly traced. State, in $200$ to
$300$ words, the chosen body, the supports, the dominant internal
force at the critical section, and the failure mode you would
inspect for first.

\paragraph{Solution sketch.} An open-ended reflection; a strong answer
names a specific member (say a lecture-hall roof beam or a parking-deck
column), draws its free body, identifies its supports (pinned, fixed,
or continuous), and states the dominant internal force at the critical
section: a bending moment $w_0L^2/8$ at midspan for a uniformly loaded
simple beam, or an axial compression for a column. The reader then
names the governing failure mode from the seven of section 8.8 (a long
roof beam: deflection or lateral-torsional buckling; a slender column:
buckling; a welded crane girder: fatigue) and what they would inspect
first, such as deflection sag, lateral bracing, or weld-toe cracking.
The value of the exercise is the habit of tracing a real load path
rather than the specific answer.
\end{exercise}
